%% 6-10 オイラーの公式（第14回：組 V6-6。ページ制限なし。偶数ページ）
%% ほかの巻が指している形：e^{jθ}=cosθ+j sinθ（5-17 のもっと知りたい。6-10 はテイラー展開以外の説明も足す）、
%%   回る矢印（フェーザ）をたすと合成になる（4-9:433）、コンデンサーの電流は電圧より 1/4 周期進む（5-15:678）
%% 3段：目でわかる（小さな回転のくり返し・回る矢印の影）→手でわかる（微分で確かめる・肩のたし算）→式でわかる
\makeatletter
% 小見出し（太字の1行）  \bkvikhead{見出し}
\newcommand\bkvikhead[1]{\par\addvspace{6pt}\Needspace{6\baselineskip}\noindent{\sffamily\bfseries\color{\bkTT}#1}\par\nobreak\vspace{2pt}}
% 囲みの中の小見出し（それだけで1行）
\newcommand\bkvikbox[1]{\par\addvspace{3pt}\noindent\textbf{#1}\par\noindent\ignorespaces}
% 大事な式のわく  \bkvikframe[はば]{中身}
\newcommand\bkvikframe[2][]{\par\addvspace{4pt}{\centering\tikz\node[draw=\bkTT, line width=0.9pt, fill=bkTheme!10, rounded corners=1.5mm, inner xsep=6mm, inner ysep=2.5mm, align=center]{#2};\par}\addvspace{4pt}\noindent\ignorespaces}
% 式を、まん中に1行で置く
\newcommand\bkvikline[1]{\par\addvspace{3pt}{\centering #1\par}\addvspace{3pt}\noindent\ignorespaces}
% 例題の途中式の行
\newcommand\bkvikq{\par\vspace{2pt}\noindent\hangindent=4.3\zw\hangafter=1\hspace*{2.6\zw}}
\newcommand\bkvika{\par\vspace{2pt}\noindent\hangindent=2.6\zw\hangafter=0\relax}
\newcommand\bkvikb{\par\vspace{2pt}\noindent\hangindent=4.3\zw\hangafter=0\relax}
% 複素平面の軸  \bkvikaxes{左}{右}{下}{上}

\newcommand\bkvikaxes[4]{%
  \draw[bkGray, line width=0.6pt, -{Stealth[length=3.5pt]}] (#1,0) -- (#2,0) node[anchor=north east, font=\sffamily\footnotesize, text=bkInk, inner sep=1pt, yshift=-1.2mm] {実部};
  \draw[bkGray, line width=0.6pt, -{Stealth[length=3.5pt]}] (0,#3) -- (0,#4) node[anchor=north west, font=\sffamily\footnotesize, text=bkInk, inner sep=1pt, xshift=1.2mm] {虚部};}
\makeatother

\bkchapter{6}{10}{オイラーの公式}{\bklevelrange{4}{5}}{オイラーの公式}
\bkmichishirube{4-8,4-9,5-14,5-15,5-17,6-9}{6-11,6-12,6-13,9-1,9-3}

\begin{goal}
\begin{itemize}[label={\color{\bkTT}●}, leftmargin=1.6\zw, itemsep=0pt, topsep=1pt]
\item \bklevel{4}\enspace $e^{j\theta}$ が「単位円の上の、角 $\theta$ の点」であることを、小さな回転のくり返しと微分で説明できる。
\item \bklevel{4}\enspace オイラーの公式で $e^{j\theta}$ を $a+bj$ に直し、$re^{j\theta}$ のかけ算・わり算・べきを「肩のたし算」で計算できる。
\item \bklevel{5}\enspace 同じ速さで回る矢印（フェーザ）のたし算で交流を合成し、「微分は $j\omega$ をかけること」を説明できる。
\end{itemize}
\end{goal}

\begin{hatena}
\noindent\begin{minipage}[c]{0.6\linewidth}
\raggedright
$e^{2}$ は「$e$ を2回かけた数」、$e^{0.5}$ は「2回かけると $e$ になる数」でした（5-14）。

では、$e^{j\theta}$ はどうでしょう。「$e$ を $j\theta$ 回かける」と言われても、何のことか分かりません。ところが電気の技術者は、交流の計算で、毎日のように $e^{j\omega t}$ を使っています。$e^{j\theta}$ とは、いったい、どんな数なのでしょう。
\end{minipage}\hfill
\begin{minipage}[c]{0.36\linewidth}
\centering
\begin{tikzpicture}[x=1mm, y=1mm]
  % 上：ふえる指数関数
  \begin{scope}[yshift=24mm]
  \draw[bkGray, line width=0.6pt, -{Stealth[length=3pt]}] (0,0) -- (34,0);
  \draw[bkGray, line width=0.6pt, -{Stealth[length=3pt]}] (0,0) -- (0,17);
  \draw[bkOrange, line width=1.4pt, domain=0:2.5, samples=40] plot ({12*\x}, {1.3*exp(\x)});
  \node[font=\sffamily\small, text=bkOrange!80!black, anchor=west, inner sep=1pt] at (3,12) {$e^{x}$};
  \end{scope}
  % 下：ゆれる三角関数
  \draw[bkGray, line width=0.6pt, -{Stealth[length=3pt]}] (0,0) -- (34,0);
  \draw[bkGray, line width=0.6pt] (0,-7.5) -- (0,7.5);
  \draw[bkBlue, line width=1.4pt, domain=0:6.2832, samples=60] plot ({5*\x}, {6*sin(\x r)});
  \node[font=\sffamily\small, text=bkBlue, anchor=west, inner sep=1pt] at (18,5.5) {$\sin x$};
  % つなぐ「？」
  \node[font=\sffamily\bfseries\Large, text=bkYellow!40!black] at (40,13) {？};
\end{tikzpicture}
\end{minipage}
\end{hatena}

\noindent 6-8 で、2乗すると $-1$ になる数 $j$ を考えました。6-9 では、複素数を平面の点（矢印）と見ると、かけ算が「回転と拡大」になることを学びました。この章では、この2つと、第5巻の指数関数 $e^x$ を組み合わせます。すると、指数関数と三角関数が1本の式でつながり、回るもの・ゆれるもの・交流が、かけ算とたし算だけで計算できるようになります。

この章は、大学で学ぶ内容です。はじめての考えが続くので、1回で分からなくても大丈夫です。「目でわかる」「手でわかる」の段だけでも、$e^{j\theta}$ の正体はつかめます。

%% ---------------- 思い出そう・やってみよう ----------------
\begin{omoidasou}
\bkrecall{4-8}{単位円の上で角 $\theta$ の点は $(\cos\theta,\ \sin\theta)$。円を回る点の影は波になる。}
\bkrecall{4-9}{加法定理：\\$\cos(a+b)=\cos a\cos b-\sin a\sin b$\\$\sin(a+b)=\sin a\cos b+\cos a\sin b$}
\bkrecall{4-9}{合成：$a\sin\theta+b\cos\theta=\sqrt{a^2+b^2}\,\sin(\theta+\alpha)$}
\bkrecall{5-14}{$\Bigl(1+\dfrac{r}{n}\Bigr)^n$ は、$n$ を大きくすると $e^r$ に近づく。$(e^{at})'=ae^{at}$}
\bkrecall{5-15}{円の上を回る点の速度は、半径に垂直。$(\sin\omega t)'=\omega\cos\omega t$。$\sin$・$\cos$ の微分はラジアンで。}
\bkrecall{5-17}{$e^x=1+x+\dfrac{x^2}{2!}+\cdots$、$\cos x=1-\dfrac{x^2}{2!}+\cdots$、$\sin x=x-\dfrac{x^3}{3!}+\cdots$}
\bkrecall{6-8}{$j^2=-1$。複素数 $a+bj$ の $a$ が実部 $\mathrm{Re}\,z$、$b$ が虚部 $\mathrm{Im}\,z$。}
\bkrecall{6-9}{複素平面。$\lvert z\rvert$ は原点からの距離、$\arg z$ は角（偏角）。極形式 $r(\cos\theta+j\sin\theta)$。\\かけ算は「大きさはかけ算、角はたし算」。$j$ をかけると90°回る。共役 $\bar z$ は横軸で裏返し。}
\end{omoidasou}

\begin{bkfigure}
\begin{tikzpicture}[x=1mm, y=1mm, nd/.style={bk node, font=\sffamily\footnotesize, text width=30mm, inner sep=1.6mm, minimum height=9mm}]
  \node[nd] (a) at (0,13) {① 小さな回転を\\くり返す（目）};
  \node[nd] (b) at (0,0) {③ 微分して\\確かめる（手）};
  \node[nd] (c) at (0,-13) {⑤ テイラー展開\\（式）};
  \node[bk node, draw=\bkTT, line width=1.3pt, fill=bkTheme!12, font=\small, text width=33mm, inner sep=2mm, minimum height=14mm] (e) at (45,0) {オイラーの公式\\$e^{j\theta}=\cos\theta+j\sin\theta$};
  \node[nd] (p) at (90,13) {⑦ かけ算は\\肩のたし算};
  \node[nd] (q) at (90,0) {⑨⑩ 交流と\\フェーザ};
  \node[nd] (r) at (90,-13) {⑪ らせん\\（6-12 へ）};
  \foreach \u in {a,b,c} {\draw[bk link] (\u.east) -- (e.west);}
  \foreach \u in {p,q,r} {\draw[bk link] (e.east) -- (\u.west);}
\end{tikzpicture}
\bkfigcaption{この章の道すじ。3つのちがう道が、同じ1本の式にたどり着く。その式で、回転と交流が計算できる。}
\end{bkfigure}

\begin{yattemiyou}[$1+0.1j$ を、10回かけてみよう]
表計算で、点 $1$ から始めて、$1+0.1j$ を何回もかけていきます。複素数 $a+bj$ は、横（実部）$a$ と、たて（虚部）$b$ の2つの列に分けて書きます。
\begin{enumerate}[label={\bkcircled[bkGreen]{\arabic*}}, leftmargin=2.4\zw, itemsep=1pt, topsep=2pt]
\item 1行めに「回数」「横（実部）」「たて（虚部）」と書く。2行めに 0、1、0 と入れる（点 $1=1+0j$）。
\item B3に \texttt{=B2-0.1*C2}、C3に \texttt{=C2+0.1*B2} と入れる（$1+0.1j$ をかける計算。わけは下）。
\item 3行めを、12行めまで下へコピーする（10回かけたことになる）。
\item 点 $(\text{横},\ \text{たて})$ を方眼紙にかき、原点から半径1の円（単位円）もかいてみよう。
\end{enumerate}
\noindent\begin{minipage}[c]{0.52\linewidth}
\centering
\bksheet{3}{15}{%
1 & 回数 & 横 & たて \\
2 & 0 & 1 & 0 \\
3 & 1 & 1 & 0.1 \\
4 & 2 & 0.99 & 0.2 \\
5 & 3 & 0.97 & 0.299 \\
6 & 4 & 0.9401 & 0.396 \\
7 & 5 & 0.9005 & 0.490 \\
8 & 6 & 0.8515 & 0.5801 \\
9 & 7 & 0.7935 & 0.6652 \\
10 & 8 & 0.7270 & 0.7446 \\
11 & 9 & 0.6525 & 0.8173 \\
12 & 10 & 0.5708 & 0.8825 \\}
\end{minipage}\hfill
\begin{minipage}[c]{0.46\linewidth}
\centering
\begin{tikzpicture}[x=46mm, y=46mm]
  \draw[bkGray, line width=0.6pt, -{Stealth[length=3.5pt]}] (-0.05,0) -- (1.16,0) node[anchor=north east, font=\sffamily\footnotesize, text=bkInk, inner sep=1pt, yshift=-1.2mm] {実部};
  \draw[bkGray, line width=0.6pt, -{Stealth[length=3.5pt]}] (0,-0.05) -- (0,1.12) node[anchor=north west, font=\sffamily\footnotesize, text=bkInk, inner sep=1pt, xshift=1.2mm] {虚部};
  \draw[bkInk!40, line width=0.8pt] (1,0) arc (0:95:1);
  \draw[bkOrange, line width=0.9pt] (1,0) -- (1,0.1) -- (0.99,0.2) -- (0.97,0.299) -- (0.9401,0.396) -- (0.9005,0.49) -- (0.8515,0.5801) -- (0.7935,0.6652) -- (0.727,0.7446) -- (0.6525,0.8173) -- (0.5708,0.8825);
  \foreach \p in {(1,0),(1,0.1),(0.99,0.2),(0.97,0.299),(0.9401,0.396),(0.9005,0.49),(0.8515,0.5801),(0.7935,0.6652),(0.727,0.7446),(0.6525,0.8173),(0.5708,0.8825)} {\fill[bkOrange!85!black] \p circle (0.45mm);}
  \fill[bkBlue] (0.5403,0.8415) circle (0.6mm);
  \draw[bkInk, line width=0.4pt] (0.5708,0.8825) -- (0.66,1.02) node[anchor=south, font=\sffamily\footnotesize, inner sep=1pt, text=bkOrange!80!black] {10回め};
  \draw[bkInk, line width=0.4pt] (0.5403,0.8415) -- (0.30,0.86) node[anchor=east, font=\sffamily\footnotesize, inner sep=1pt, align=right, text=bkBlue] {角1の点};
  \node[font=\sffamily\footnotesize, text=bkGray, anchor=north east, inner sep=1pt] at (0.62,0.62) {単位円};
\end{tikzpicture}
\end{minipage}
\end{yattemiyou}

\noindent 点は、ほとんど単位円の上を、左回りに少しずつ進みました。10回めの点 $(0.571,\ 0.883)$ は、単位円の上で角1ラジアンの所にある点（青）の近くです。

\bkvikhead{B3・C3の式のわけ}
今の点を $a+bj$（B2が $a$、C2が $b$）とします。$1+0.1j$ をかけて、かっこを外します（3-8 の展開、6-8 の $j^2=-1$）。
\begin{bkhenkei}
(a+bj)(1+0.1j) & かけ算の式 \\
=a+0.1aj+bj+0.1bj^2 & かっこを外す（4つの項） \\
=a+0.1aj+bj-0.1b & $j^2=-1$ を使う \\
=(a-0.1b)+(b+0.1a)j & 実部と虚部に分ける \\
\end{bkhenkei}
新しい横は $a-0.1b$（B3の式）、新しいたては $b+0.1a$（C3の式）です。

\noindent なぜ点は円の上を回るのでしょう。そして、10回めの点は、なぜ角1の点の近くに来たのでしょう。これが、この章の入口です。



%% ---------------- 数学の道具 ----------------
\bksection{doguu}{数学の道具}

\bkstage{me}{目でわかる}{小さな回転を、何回もくり返す}{\bklevel{4}}

\bksubsection{1}{$1+0.1j$ をかけると、何が起きる？}

\label{bk:6-10:kaiten}6-9 で学んだように、複素数のかけ算は「大きさはかけ算、角はたし算」です。そこで、かける数 $1+0.1j$ の大きさと角を調べます。

\begin{bkfigure}
\begin{tikzpicture}[x=1mm, y=1mm]
  \draw[bkGray, line width=0.6pt, -{Stealth[length=3.5pt]}] (-4,0) -- (92,0) node[anchor=north east, font=\sffamily\footnotesize, text=bkInk, inner sep=1pt, yshift=-1.2mm] {実部};
  \draw[bkGray, line width=0.6pt, -{Stealth[length=3.5pt]}] (0,-4) -- (0,16) node[anchor=north west, font=\sffamily\footnotesize, text=bkInk, inner sep=1pt, xshift=1.2mm] {虚部};
  \draw[bkBlue, line width=1.6pt] (0,0) -- (80,0);
  \draw[bkOrange, line width=1.6pt] (80,0) -- (80,8);
  \draw[bkInk, line width=1.1pt] (0,0) -- (80,8);
  \fill[bkInk] (80,8) circle (1);
  \node[font=\sffamily\footnotesize, anchor=north, text=bkBlue, inner sep=1pt] at (62,-1.5) {横 1};
  \node[font=\sffamily\footnotesize, anchor=west, text=bkOrange!80!black, inner sep=1pt] at (81.5,3.5) {たて 0.1};
  \node[font=\sffamily\footnotesize, anchor=south, inner sep=1pt] at (80,9.5) {$1+0.1j$};
  \node[font=\sffamily\footnotesize, anchor=south east, inner sep=1pt] at (48,5.6) {長さ $\sqrt{1.01}=1.005$};
  \draw[bkInk, line width=0.5pt] (24,0) arc (0:5.71:24);
  \draw[bkInk, line width=0.4pt] (24.3,1.2) -- (21,-3.2) node[anchor=north, font=\sffamily\footnotesize, inner sep=1pt] {角 約0.1ラジアン};
\end{tikzpicture}
\bkfigcaption{$1+0.1j$ は、長さがほとんど1で、角が約0.1ラジアンの矢印。}
\end{bkfigure}
\begin{itemize}[leftmargin=1.2\zw, itemsep=2pt, topsep=2pt, label={\color{\bkTT}●}]
\item \textbf{大きさ}\quad 横1、たて0.1の直角三角形の斜辺です（4-4 の三平方）。長さは $\sqrt{1^2+0.1^2}=\sqrt{1.01}$ で、$1.00499\cdots$。ほとんど1です。
\item \textbf{角}\quad $\tan^{-1}0.1=0.0997$ ラジアンです（4-15）。ほとんど0.1ラジアンです。
\end{itemize}
だから、$1+0.1j$ をかけるたびに、点は「長さはほとんど変わらず、角が約0.1ラジアンだけ左へ回る」のです。10回かけると、角は約 $0.1\times10=1$ ラジアン、長さは $1.005^{10}=1.051$ 倍になります。だから、10回めの点は、単位円の上の角1の点の、少し外側に来ました。

\bkvikhead{もっと細かく回すと：$\Bigl(1+\dfrac{j}{n}\Bigr)^n$}
目標の「角1ラジアン」を、$n$ 回に分けて回すことにします。1回分は $1+\dfrac{j}{n}$ をかけます（$n=10$ が、やってみようの $1+0.1j$）。$n$ をふやすと、点はどうなるでしょう。
\begin{bkcenter}
\begin{bktable}{colspec={Q[c,m,wd=16mm] Q[c,m,wd=34mm] Q[c,m,wd=22mm] Q[c,m,wd=26mm]}, rowsep=1.5pt}
$n$ & $\Bigl(1+\dfrac{j}{n}\Bigr)^n$ & 大きさ & 角（ラジアン） \\
1 & $1+j$ & 1.414 & 0.785 \\
2 & $0.75+j$ & 1.25 & 0.927 \\
4 & $0.629+0.938j$ & 1.129 & 0.980 \\
10 & $0.571+0.883j$ & 1.051 & 0.997 \\
100 & $0.543+0.846j$ & 1.005 & 0.99997 \\
\num{10000} & $0.5403+0.8415j$ & 1.00005 & 1.00000 \\
\end{bktable}
\bktabcaption{$n$ をふやすと、大きさは1に、角は1ラジアンに近づく。}
\end{bkcenter}
$n=2$ なら、$\Bigl(1+\dfrac{j}{2}\Bigr)^2=1+j+\dfrac{j^2}{4}=1+j-0.25=0.75+j$ です（手で計算できます）。$n$ を大きくすると、点は $\cos1+j\sin1=0.5403+0.8415j$、つまり単位円の上の角1の点に近づいていきます。

\begin{bkfigure}
\begin{tikzpicture}[x=56mm, y=56mm]
  \draw[bkGray, line width=0.6pt, -{Stealth[length=3.5pt]}] (-0.05,0) -- (1.2,0) node[anchor=north east, font=\sffamily\footnotesize, text=bkInk, inner sep=1pt, yshift=-1.2mm] {実部};
  \draw[bkGray, line width=0.6pt, -{Stealth[length=3.5pt]}] (0,-0.05) -- (0,1.14) node[anchor=north west, font=\sffamily\footnotesize, text=bkInk, inner sep=1pt, xshift=1.2mm] {虚部};
  \draw[bkInk!40, line width=0.8pt] (1,0) arc (0:90:1);
  \draw[bkGray, line width=1.1pt] (1,0) -- (1,0.5) -- (0.75,1);
  \draw[bkPurple, line width=1.1pt] (1,0) -- (1,0.25) -- (0.9375,0.5) -- (0.8125,0.734375) -- (0.62891,0.9375);
  \draw[bkOrange, line width=1.1pt] (1,0) -- (1,0.1) -- (0.99,0.2) -- (0.97,0.299) -- (0.9401,0.396) -- (0.9005,0.49) -- (0.8515,0.5801) -- (0.7935,0.6652) -- (0.727,0.7446) -- (0.6525,0.8173) -- (0.5708,0.8825);
  \fill[bkGray] (0.75,1) circle (0.7mm);
  \fill[bkPurple] (0.62891,0.9375) circle (0.7mm);
  \fill[bkOrange!85!black] (0.5708,0.8825) circle (0.7mm);
  \fill[bkBlue] (0.5403,0.8415) circle (0.8mm);
  \draw[bkInk, line width=0.4pt] (0,0) -- (0.5403,0.8415);
  \draw[bkInk, line width=0.5pt] (0.16,0) arc (0:57.3:0.16);
  \node[font=\sffamily\footnotesize] at (28.6:0.21) {1};
  % 線の見本（右下のあいている所）
  \draw[bkGray, line width=1.1pt] (0.42,0.30) -- (0.52,0.30) node[anchor=west, font=\sffamily\footnotesize, text=bkInk, inner sep=1.5pt] {$n=2$};
  \draw[bkPurple, line width=1.1pt] (0.42,0.22) -- (0.52,0.22) node[anchor=west, font=\sffamily\footnotesize, text=bkInk, inner sep=1.5pt] {$n=4$};
  \draw[bkOrange, line width=1.1pt] (0.42,0.14) -- (0.52,0.14) node[anchor=west, font=\sffamily\footnotesize, text=bkInk, inner sep=1.5pt] {$n=10$};
  \draw[bkInk, line width=0.4pt] (0.5403,0.8415) -- (0.49,0.935) node[anchor=south east, font=\sffamily\footnotesize, text=bkBlue, inner sep=1pt] {角1の点};
  \node[font=\sffamily\footnotesize, text=bkGray, anchor=north east, inner sep=1pt] at (0.70,0.66) {単位円};
\end{tikzpicture}
\bkfigcaption{1回分の回転を小さく（$n$ を大きく）するほど、折れ線は単位円にそい、終わりの点は角1の点（青い点 $\cos1+j\sin1$）に近づく。}
\end{bkfigure}

\bkvikhead{角 $\theta$ でも同じ}
1ラジアンでなく角 $\theta$ だけ回したいときは、1回分を $1+\dfrac{j\theta}{n}$ にします。この数の角は、ほとんど $\dfrac{\theta}{n}$ です（小さい角では、$\tan^{-1}x\approx x$。5-17 の1次近似）。大きさは $\sqrt{1+\dfrac{\theta^2}{n^2}}$ で、$n$ が大きいとほとんど1です。だから $n$ 回かけると、角は $\dfrac{\theta}{n}\times n=\theta$ に、大きさは1に近づきます。
\bkvikframe{$\Bigl(1+\dfrac{j\theta}{n}\Bigr)^n$ は、$n$ を大きくすると、\\単位円の上の角 $\theta$ の点 $\cos\theta+j\sin\theta$ に近づく}

\bkvikhead{\bkhatenaref の答え}
ここで、5-14 を思い出しましょう。$\Bigl(1+\dfrac{r}{n}\Bigr)^n$ は、$n$ を大きくすると $e^r$ に近づきました。上の式は、この $r$ のところに $j\theta$ を入れたものです。そこで、この行き着く先を $e^{j\theta}$ と書くことにします。
\bkvikframe{\Large $e^{j\theta}=\cos\theta+j\sin\theta$}
これを\textbf{オイラーの公式}といいます。\bkidx{おいらーのこうしき}{オイラーの公式}$e^{j\theta}$ は「$e$ を $j\theta$ 回かけた数」ではありません。「$1+\dfrac{j\theta}{n}$ という、ほんの少しだけ回す数を、何回もかけた行き着く先」で、それは\textbf{単位円の上の、角 $\theta$ の点}です。

\begin{seikaku}
5-14 までの $e^x$ は、$x$ が実数のときの関数でした。$x$ が複素数のときの $e^x$ は、まだ決めていなかったので、ここで $e^{j\theta}=\cos\theta+j\sin\theta$ と\textbf{決めた}（定義した）のです。好きに決めてよいなら、なぜこう決めるのでしょう。それは、こう決めると、指数関数のきまり（$e^ae^b=e^{a+b}$、$(e^{at})'=ae^{at}$、テイラー展開）が、そのまま成り立つからです。このあとの③④⑤で、1つずつ確かめます。
\end{seikaku}

\bksubsection{2}{回る矢印と、その影}

\label{bk:6-10:kage}$\theta$ のかわりに $\omega t$（$\omega$ は角速度〔rad/s〕、4-8）を入れると、
\bkvikline{$e^{j\omega t}=\cos\omega t+j\sin\omega t$}
\begin{bkfigure}
\begin{tikzpicture}[x=0.82mm, y=0.82mm]
  % 円
  \draw[bkGray, line width=0.6pt] (-18,0) -- (18,0);
  \draw[bkGray, line width=0.6pt] (0,-18) -- (0,18);
  \draw[bkInk!45, line width=0.8pt] (0,0) circle (14);
  \draw[bkBlue, line width=1.4pt, -{Stealth[length=4.5pt]}] (0,0) -- (50:14);
  \fill[bkBlue] (50:14) circle (1);
  \draw[bkInk, line width=0.5pt] (4.5,0) arc (0:50:4.5);
  \node[font=\sffamily\footnotesize] at (25:7.4) {$\omega t$};
  \node[font=\sffamily\footnotesize, text=bkBlue, anchor=south west, inner sep=1pt] at ($(50:14)+(1.4,0.9)$) {$e^{j\omega t}$};
  % 右：たての影（sin）
  \draw[bkGray, line width=0.6pt, -{Stealth[length=3.5pt]}] (30,0) -- (96,0) node[anchor=west, font=\sffamily\footnotesize, text=bkInk] {$\omega t$};
  \draw[bkGray, line width=0.6pt] (30,-16) -- (30,16);
  \draw[bkOrange, line width=1.5pt, domain=0:6.2832, samples=80] plot ({30+9.549*\x}, {14*sin(\x r)});
  \draw[bkGray, dashed, line width=0.5pt] (50:14) -- (38.33,10.72);
  \fill[bkOrange!85!black] (38.33,10.72) circle (1);
  \node[font=\sffamily\footnotesize, anchor=west, text=bkOrange!80!black, fill=white, inner sep=1pt] at (60,13) {たての影（虚部）$\sin\omega t$};
  % 下：横の影（cos）
  \draw[bkGray, line width=0.6pt, -{Stealth[length=3.5pt]}] (0,-22) -- (0,-66) node[anchor=north, font=\sffamily\footnotesize, text=bkInk] {$\omega t$};
  \draw[bkGray, line width=0.6pt] (-16,-22) -- (16,-22);
  \draw[bkPurple, line width=1.5pt, domain=0:6.2832, samples=80] plot ({14*cos(\x r)}, {-22-6.366*\x});
  \draw[bkGray, dashed, line width=0.5pt] (50:14) -- (9.0,-27.56);
  \fill[bkPurple] (9.0,-27.56) circle (1);
  \node[font=\sffamily\footnotesize, anchor=west, text=bkPurple, fill=white, inner sep=1pt, align=left] at (18,-44) {横の影（実部）\\$\cos\omega t$};
\end{tikzpicture}
\bkfigcaption{回る矢印 $e^{j\omega t}$（青）。たての位置（虚部）は $\sin$ の波（右）、横の位置（実部）は $\cos$ の波（下。時間は下へ進む）になる。}
\end{bkfigure}
$t$ がふえると、点 $e^{j\omega t}$ は、単位円の上を1秒に $\omega$ ラジアンずつ、左回りに回ります。4-8 で、回る点の影が波になることを学びました。ここでも同じです。
\begin{itemize}[leftmargin=1.2\zw, itemsep=2pt, topsep=2pt, label={\color{\bkTT}●}]
\item 横の位置（実部）は $\cos\omega t$。真上から光を当てたときの、床の影です。
\item たての位置（虚部）は $\sin\omega t$。真横から光を当てたときの、かべの影です。
\end{itemize}


1つの $e^{j\omega t}$ の中に、$\cos$ の波と $\sin$ の波が、両方入っているのです。矢印をあつかうほうが、波を2本あつかうより楽になることが多い。これが、電気の技術者が $e^{j\omega t}$ を使うわけです（⑨⑩でくわしく）。

%% ---------------- 手でわかる ----------------
\Needspace{11\baselineskip}
\bkstage{te}{手でわかる}{微分して確かめる}{\bklevel{4}}

\bksubsection{3}{$e^{jt}$ を微分すると：速度は、位置を90°回したもの}

\label{bk:6-10:bibun}①とは別の道から、オイラーの公式にたどり着いてみます。手がかりは、5-14 の微分のきまり $(e^{at})'=ae^{at}$ です。このきまりが $a=j$ でも成り立つとしたら、$z=e^{jt}$ はどんな動きをするでしょう。

\begin{enumerate}[label={\bkcircled[\bkTF]{\arabic*}}, leftmargin=2.4\zw, itemsep=2pt, topsep=2pt]
\item $a=j$ を入れると、$\dfrac{\dd z}{\dd t}=je^{jt}=jz$。$z$ を位置と見ると、$\dfrac{\dd z}{\dd t}$ は速度です（5-5）。だから「\textbf{速度＝位置に $j$ をかけたもの}」。
\item $j$ をかけると、矢印は長さを変えずに、左回りに90°回ります（6-9）。だから速度の矢印は、位置の矢印と\textbf{同じ長さで、90°回った向き}です。
\item はじめ（$t=0$）は $e^0=1$ なので、点は1から出発します。位置の矢印の長さは1、速度の長さも1です。
\item 速度がいつも半径に垂直なので、点は原点からの距離を変えずに、円をえがいて回ります（5-15①：円の上を回る点の速度は半径に垂直）。速さは1なので、$t$ 秒で進む弧の長さは $t$、回る角は $t$ ラジアンです（弧の長さ＝半径×角、4-2）。
\end{enumerate}
だから、$e^{jt}$ は「単位円の上の角 $t$ の点」、つまり $\cos t+j\sin t$ です。①と同じ答えにたどり着きました。

\begin{bkfigure}
\begin{tikzpicture}[x=1mm, y=1mm]
  \begin{scope}
  \bkvikaxes{-6}{30}{-6}{31}
  \draw[bkInk!45, line width=0.8pt] (0,0) circle (16);
  \draw[bkBlue, line width=1.4pt, -{Stealth[length=5pt]}] (0,0) -- (35:16);
  \draw[bkOrange, line width=1.4pt, -{Stealth[length=5pt]}] (35:16) -- ++(125:16);
  \draw[bkInk, line width=0.4pt] ($(35:16)+(215:2.6)$) -- ++(125:2.6) -- ++(35:2.6);
  \node[font=\sffamily\footnotesize, text=bkBlue, anchor=south east, inner sep=1pt] at (7.2,6.6) {$e^{jt}$};
  \node[font=\sffamily\footnotesize, text=bkOrange!80!black, anchor=west, inner sep=1pt, align=left] at (11,21) {速度 $je^{jt}$\\（90°回した）};
  \end{scope}
  \begin{scope}[xshift=68mm]
  \bkvikaxes{-6}{30}{-6}{31}
  \draw[bkInk!45, line width=0.8pt] (0,0) circle (20);
  \draw[bkOrange, line width=1.1pt] (20,0) -- (20,7) -- (17.55,14) -- (12.65,20.14) -- (5.6,24.57);
  \foreach \p in {(20,0),(20,7),(17.55,14),(12.65,20.14),(5.6,24.57)} {\fill[bkOrange!85!black] \p circle (0.8);}
  \node[font=\sffamily\footnotesize, anchor=west, inner sep=1pt, align=left] at (23,14) {接線の向きに\\まっすぐ進むと\\少し外へ出る};
  \end{scope}
\end{tikzpicture}
\bkfigcaption{左：速度（オレンジ）は、位置（青）を90°回した向き。右：やってみようの計算は、「今の速度で、少しの時間まっすぐ進む」をくり返したもの（6-2 のオイラー法）。1歩ごとに円の少し外へ出るので、点は少しずつふくらんだ。}
\end{bkfigure}

\bkvikhead{やってみようは、オイラー法だった}
やってみようの「$1+0.1j$ をかける」は、
\bkvikline{次の位置 $=z+jz\times0.1$\qquad（今の位置＋今の速度×0.1）}
と同じです（$z(1+0.1j)=z+0.1jz$）。これは、$\dfrac{\dd z}{\dd t}=jz$ を、6-2 のオイラー法で $\Delta t=0.1$ ずつ進めた計算です。6-2 で、ふつうのオイラー法でばねを動かすと、ゆれがふくらむことを見ました。ここで点が少しずつ外へ出たのも、同じ理由です。$\Delta t$ を小さく（$n$ を大きく）すると、ふくらみは小さくなり、点は円にそって回るようになります（①の表）。

\bkvikhead{逆向きに確かめる}
$\cos t+j\sin t$ を微分して、本当に $j$ 倍になるかを確かめます（5-15 の $\sin$・$\cos$ の微分）。
\begin{bkhenkei}
(\cos t+j\sin t)' & 微分する（$j$ は定数） \\
=-\sin t+j\cos t & $(\cos t)'=-\sin t$、$(\sin t)'=\cos t$ \\
\end{bkhenkei}
一方、$j$ をかけると、
\begin{bkhenkei}
j(\cos t+j\sin t) & $j$ をかける \\
=j\cos t+j^2\sin t & かっこを外す \\
=-\sin t+j\cos t & $j^2=-1$ を使い、実部を前に \\
\end{bkhenkei}
2つは同じです。$\cos t+j\sin t$ は、たしかに「微分すると自分の $j$ 倍」で、$t=0$ で1になります。5-14 のきまりが、$e^{jt}$ でもそのまま成り立っています。

\Needspace{12\baselineskip}
\bksubsection{4}{かけ算は「肩のたし算」}

\label{bk:6-10:kakezan}指数関数のいちばん大切なきまりは、$e^ae^b=e^{a+b}$（かけ算が、肩のたし算になる）でした（3-11・5-14）。$e^{j\theta}$ でも成り立つでしょうか。

\noindent\begin{minipage}[c]{0.55\linewidth}
\raggedright
\textbf{図で}\quad $e^{ja}$ は角 $a$、$e^{jb}$ は角 $b$ の、長さ1の矢印です。6-9 のきまり「かけ算は大きさをかけ、角をたす」から、$e^{ja}\times e^{jb}$ は、長さ $1\times1=1$、角 $a+b$ の矢印。これは $e^{j(a+b)}$ です。
\[ e^{ja}\,e^{jb}=e^{j(a+b)} \]
\end{minipage}\hfill
\begin{minipage}[c]{0.4\linewidth}
\centering
\begin{tikzpicture}[x=1mm, y=1mm]
  \bkvikaxes{-4}{30}{-4}{29}
  \draw[bkInk!45, line width=0.8pt] (22,0) arc (0:95:22);
  \draw[bkBlue, line width=1.3pt, -{Stealth[length=4.5pt]}] (0,0) -- (20:22);
  \draw[bkOrange, line width=1.3pt, -{Stealth[length=4.5pt]}] (0,0) -- (35:22);
  \draw[bkPurple, line width=1.5pt, -{Stealth[length=4.5pt]}] (0,0) -- (55:22);
  \node[font=\sffamily\footnotesize, text=bkBlue, anchor=west, inner sep=1pt] at ($(20:22)+(0.8,0)$) {$e^{ja}$};
  \node[font=\sffamily\footnotesize, text=bkOrange!80!black, anchor=south west, inner sep=1pt] at ($(35:22)+(0.6,0.2)$) {$e^{jb}$};
  \node[font=\sffamily\footnotesize, text=bkPurple, anchor=south west, inner sep=1pt] at ($(55:22)+(0.4,0.4)$) {$e^{j(a+b)}$};
  \draw[bkInk, line width=0.5pt] (9,0) arc (0:20:9);
  \node[font=\sffamily\footnotesize] at (10:11.6) {$a$};
\end{tikzpicture}
\end{minipage}

\bkvikhead{式で}
左のかけ算を、オイラーの公式で書いて広げます。
\begin{bkhenkei}[0.75]
e^{ja}e^{jb}=(\cos a+j\sin a)(\cos b+j\sin b) & オイラーの公式 \\
=\cos a\cos b+j\cos a\sin b+j\sin a\cos b+j^2\sin a\sin b & 展開する \\
=(\cos a\cos b-\sin a\sin b)+j(\sin a\cos b+\cos a\sin b) & $j^2=-1$、実部と虚部に分ける \\
=\cos(a+b)+j\sin(a+b) & 4-9 の加法定理 \\
=e^{j(a+b)} & オイラーの公式 \\
\end{bkhenkei}
実部が $\cos$ の加法定理、虚部が $\sin$ の加法定理になりました。4-9 で苦労して作った2つの加法定理が、「肩のたし算」という1つのきまりにまとまっています。逆に言えば、加法定理を忘れても、$e^{ja}e^{jb}=e^{j(a+b)}$ を広げれば、いつでも作れます。

\bkvikhead{マイナスの角：$e^{-j\theta}$}
$\theta$ に $-\theta$ を入れると、$\cos(-\theta)=\cos\theta$、$\sin(-\theta)=-\sin\theta$（4-8）なので、
\bkvikline{$e^{-j\theta}=\cos\theta-j\sin\theta$}
これは $e^{j\theta}$ を横軸で裏返した点、つまり共役 $\overline{e^{j\theta}}$ です（6-9）。右回りに $\theta$ だけ回した点です。また、肩のたし算から $e^{j\theta}e^{-j\theta}=e^{0}=1$ なので、$e^{-j\theta}=\dfrac{1}{e^{j\theta}}$ にもなっています。

\bkkokomade[$e^{j\theta}$ は単位円の上の角 $\theta$ の点。かけると $\theta$ 回る。]

%% ---------------- 式でわかる ----------------
\Needspace{11\baselineskip}
\bkstage{shiki}{式でわかる}{テイラー展開・極形式・交流}{\bklevel{5}}

\bksubsection{5}{3つめの道：テイラー展開}

\label{bk:6-10:taylor}5-17 のもっと知りたいでは、$e^x$ のテイラー展開の $x$ に $j\theta$ を入れて、オイラーの公式を作りました。もう一度、1行ずつ見てみましょう。まず、$j$ のべきは4つごとにくり返します（6-8）。
\begin{bkcenter}
\begin{tblr}{colspec={Q[c,m,wd=10mm]*{8}{Q[c,m,wd=9mm]}}, row{1}={bg=bkTheme!12, font=\sffamily\small}, hline{1,2,Z}={0.6pt,bkGray}, colsep=2pt}
$n$ & 0 & 1 & 2 & 3 & 4 & 5 & 6 & 7 \\
$j^n$ & 1 & $j$ & $-1$ & $-j$ & 1 & $j$ & $-1$ & $-j$ \\
\end{tblr}
\end{bkcenter}
$e^x=1+x+\dfrac{x^2}{2!}+\dfrac{x^3}{3!}+\cdots$ に $x=j\theta$ を入れます。$(j\theta)^n=j^n\theta^n$ なので、表の $j^n$ を使って、
\[ e^{j\theta}=1+j\theta-\frac{\theta^2}{2!}-j\frac{\theta^3}{3!}+\frac{\theta^4}{4!}+j\frac{\theta^5}{5!}-\cdots \]
$j$ のつかない項と、$j$ のつく項に分けると、
\[ e^{j\theta}=\Bigl(1-\frac{\theta^2}{2!}+\frac{\theta^4}{4!}-\cdots\Bigr)+j\Bigl(\theta-\frac{\theta^3}{3!}+\frac{\theta^5}{5!}-\cdots\Bigr) \]
前のかっこは $\cos\theta$、後のかっこは $\sin\theta$ の展開そのものです（5-17⑤）。

\bkvikhead{数で確かめる（$\theta=1$）}
項を1つずつたしていくと、次のようになります（小数第4位まで）。
\begin{bkcenter}
\begin{bktable}{colspec={Q[l,m,wd=22mm] *{7}{Q[c,m,wd=12.5mm]}}, rowsep=1.5pt, colsep=1.5pt, cells={font=\small}}
たした項 & 1 & $j$ & $-\dfrac{1}{2}$ & $-\dfrac{j}{6}$ & $\dfrac{1}{24}$ & $\dfrac{j}{120}$ & $-\dfrac{1}{720}$ \\
実部の合計 & 1 & 1 & 0.5 & 0.5 & 0.5417 & 0.5417 & 0.5403 \\
虚部の合計 & 0 & 1 & 1 & 0.8333 & 0.8333 & 0.8417 & 0.8417 \\
\end{bktable}
\bktabcaption{$\cos1=0.5403$、$\sin1=0.8415$ に近づく。①の表の行き着く先と同じ点。}
\end{bkcenter}
①（小さな回転のくり返し）、③（微分）、⑤（テイラー展開）。3つのちがう道が、すべて同じ $\cos\theta+j\sin\theta$ にたどり着きました。$e^{j\theta}$ をこう決めることが、とても自然だと分かります。

\bksubsection{6}{特別な角と、$e^{j\pi}=-1$}

\label{bk:6-10:tokubetsu}よく使う角の $e^{j\theta}$ を、表にします。単位円の上の点を読めば、すぐに分かります。
\begin{bkcenter}
\begin{bktable}{colspec={Q[c,m,wd=11mm] Q[c,m,wd=11mm] Q[c,m,wd=24mm] Q[c,m,wd=11mm] Q[c,m,wd=11mm] Q[c,m,wd=24mm]}, rowsep=2pt, colsep=2pt}
$\theta$ & （度） & $e^{j\theta}$ & $\theta$ & （度） & $e^{j\theta}$ \\
0 & 0° & 1 & $\dfrac{\pi}{2}$ & 90° & $j$ \\
$\dfrac{\pi}{6}$ & 30° & $\dfrac{\sqrt3}{2}+\dfrac{1}{2}j$ & $\pi$ & 180° & $-1$ \\
$\dfrac{\pi}{3}$ & 60° & $\dfrac{1}{2}+\dfrac{\sqrt3}{2}j$ & $\dfrac{3\pi}{2}$ & 270° & $-j$ \\
\SetCell[c=3]{c}{} & & & $2\pi$ & 360° & 1 \\
\end{bktable}
\end{bkcenter}
\begin{itemize}[leftmargin=1.2\zw, itemsep=2pt, topsep=2pt, label={\color{\bkTT}●}]
\item $e^{j\pi/2}=j$：90°回す数は $j$（6-8 で $j$ を「90°回す数」として作ったことと合う）。
\item $e^{j\pi}=-1$：180°回すと反対向き（2-2 の「×($-1$) で向きが反対」）。
\item $e^{j2\pi}=1$：1周回ると、もとにもどる。だから $e^{j(\theta+2\pi)}=e^{j\theta}e^{j2\pi}=e^{j\theta}$。$e^{j\theta}$ は、周期 $2\pi$ でくり返します。
\end{itemize}

\noindent とくに $e^{j\pi}=-1$ を書きかえた
\bkvikline{$e^{j\pi}+1=0$}
は\textbf{オイラーの等式}とよばれます。\bkidx{おいらーのとうしき}{オイラーの等式}自然対数の底 $e$、円周率 $\pi$、2乗して $-1$ の $j$、そして1と0。べつべつの所から生まれた数が、1本の式にならんでいます。オイラーは、1748年の本『無限解析入門』で、もとになる公式を書きました。

\bkvikhead{Python で確かめる}
Python では、$j$ を \texttt{1j} と書きます（6-8）。\texttt{cmath} という道具で、複素数の $e^z$ を計算できます。
\begin{bkcode}[Python]
import cmath, math
print(cmath.exp(1j*math.pi/3))   \# e の j×π/3 乗
print(cmath.exp(1j*math.pi))     \# e の j×π 乗
\end{bkcode}
\begin{bkcodeout}
(0.5000000000000001+0.8660254037844386j)
(-1+1.2246467991473532e-16j)
\end{bkcodeout}
1行めは $\dfrac{1}{2}+\dfrac{\sqrt3}{2}j$（$\dfrac{\sqrt3}{2}=0.866\cdots$）です。2行めの \texttt{1.22e-16} は $1.22\times10^{-16}$ のことで、コンピュータの丸めのずれ（2-19）です。本当は0なので、$e^{j\pi}=-1$ です。

\bksubsection{7}{極形式 $re^{j\theta}$：かけ算・わり算・べき}

\label{bk:6-10:kyoku}6-9 では、複素数を極形式 $z=r(\cos\theta+j\sin\theta)$ で書きました（$r=\lvert z\rvert$、$\theta=\arg z$）。かっこの中はオイラーの公式そのものなので、
\bkvikframe{$z=r(\cos\theta+j\sin\theta)=re^{j\theta}$}
と短く書けます。この書き方だと、6-9 のきまりが、指数法則（3-11）としてそのまま出てきます。
\begin{bkcenter}
\begin{bktable}{colspec={Q[l,m,wd=13mm] Q[l,m,wd=61mm] Q[l,m,wd=33mm]}, rowsep=2pt}
計算 & $re^{j\theta}$ で書くと & 読み方 \\
かけ算 & $r_1e^{j\theta_1}\times r_2e^{j\theta_2}=r_1r_2\,e^{j(\theta_1+\theta_2)}$ & 大きさはかけ算、角はたし算 \\
わり算 & $\dfrac{r_1e^{j\theta_1}}{r_2e^{j\theta_2}}=\dfrac{r_1}{r_2}\,e^{j(\theta_1-\theta_2)}$ & 大きさはわり算、角はひき算 \\
べき & $\bigl(re^{j\theta}\bigr)^n=r^n\,e^{jn\theta}$ & 大きさは $n$ 乗、角は $n$ 倍 \\
\end{bktable}
\bktabcaption{6-9 のきまりは、$e^{a}e^{b}=e^{a+b}$、$\dfrac{e^a}{e^b}=e^{a-b}$、$(e^a)^n=e^{na}$ と同じ形。}\label{bk:6-10:tab-kyoku}
\end{bkcenter}
べきの行で $r=1$ にしたものが、6-9 のド・モアブルの定理です。
\bkvikline{$(\cos\theta+j\sin\theta)^n=\cos n\theta+j\sin n\theta$}

\bkvikhead{例：$(1+j)^8$}
$1+j$ は、大きさ $\sqrt{1^2+1^2}=\sqrt2$、角 $\dfrac{\pi}{4}$（45°）なので、$1+j=\sqrt2\,e^{j\pi/4}$。
\begin{bkhenkei}
(1+j)^8=\bigl(\sqrt2\,e^{j\pi/4}\bigr)^8 & 極形式に直す \\
=(\sqrt2)^8\,e^{j\,8\times\pi/4} & 大きさは8乗、角は8倍 \\
=16\,e^{j2\pi} & $(\sqrt2)^8=2^4=16$ \\
=16 & $e^{j2\pi}=1$（1周） \\
\end{bkhenkei}
$(1+j)$ を8回かけて展開するより、ずっと楽です。

\bksubsection{8}{$\cos$ と $\sin$ を、$e^{j\theta}$ で書く}

\label{bk:6-10:cossin}$e^{j\theta}=\cos\theta+j\sin\theta$ と、④の $e^{-j\theta}=\cos\theta-j\sin\theta$ を、たしたり、ひいたりします。

\noindent\begin{minipage}[c]{0.56\linewidth}
\raggedright
\textbf{たす}\quad $j\sin\theta$ が打ち消しあって、
\[ e^{j\theta}+e^{-j\theta}=2\cos\theta \]
\textbf{ひく}\quad $\cos\theta$ が打ち消しあって、
\[ e^{j\theta}-e^{-j\theta}=2j\sin\theta \]
それぞれ2と $2j$ でわると、
\[ \cos\theta=\frac{e^{j\theta}+e^{-j\theta}}{2},\quad \sin\theta=\frac{e^{j\theta}-e^{-j\theta}}{2j} \]
\end{minipage}\hfill
\begin{minipage}[c]{0.4\linewidth}
\centering
\begin{tikzpicture}[x=1mm, y=1mm]
  \bkvikaxes{-4}{40}{-21}{21}
  \draw[bkInk!45, line width=0.7pt] (40:20) arc (40:-40:20);
  \draw[bkBlue, line width=1.3pt, -{Stealth[length=4.5pt]}] (0,0) -- (40:20);
  \draw[bkOrange, line width=1.3pt, -{Stealth[length=4.5pt]}] (0,0) -- (-40:20);
  \draw[bkGray, dashed, line width=0.5pt] (40:20) -- (30.64,0) -- (-40:20);
  \draw[bkPurple, line width=1.8pt, -{Stealth[length=5pt]}] (0,0) -- (30.64,0);
  \node[font=\sffamily\footnotesize, text=bkBlue, anchor=south west, inner sep=1pt] at ($(40:20)+(0.4,0.2)$) {$e^{j\theta}$};
  \node[font=\sffamily\footnotesize, text=bkOrange!80!black, anchor=north west, inner sep=1pt] at ($(-40:20)+(0.4,-0.2)$) {$e^{-j\theta}$};
  \node[font=\sffamily\footnotesize, text=bkPurple, anchor=south, inner sep=1pt] at (19,1.3) {$2\cos\theta$};
\end{tikzpicture}
\end{minipage}

\noindent 図で見ると、左回りに回る矢印 $e^{j\theta}$ と、右回りに回る矢印 $e^{-j\theta}$ は、いつも横軸について対称です。2本をたすと、たての成分が打ち消しあい、和は横軸の上だけを行ったり来たりします。「\textbf{反対向きに回る2本の矢印をたすと、まっすぐゆれる}」。$\cos\theta$ の波は、2つの回転の和だったのです。この見方は、9-3 のフーリエ級数で大活躍します。また、6-11 の双曲線関数は、この式の $j\theta$ を $x$ にかえた形から生まれます。

\bkvikhead{$\theta=\dfrac{\pi}{3}$ で確かめる}
⑥の表から $e^{j\pi/3}=\dfrac{1}{2}+\dfrac{\sqrt3}{2}j$、$e^{-j\pi/3}=\dfrac{1}{2}-\dfrac{\sqrt3}{2}j$（共役）。たすと1なので、$\cos\dfrac{\pi}{3}=\dfrac{1}{2}$。ひくと $\sqrt3\,j$ なので、$\sin\dfrac{\pi}{3}=\dfrac{\sqrt3\,j}{2j}=\dfrac{\sqrt3}{2}$。どちらも合っています。


% Ver2（No.34）：3倍角の「もっと知りたい」を⑧のあとへ移した（⑦の例のあとのページの下の空きをへらす）
\Needspace{22\baselineskip}
\begin{motto}[5]
\bkvikbox{3倍角の式を、オイラーの公式で作る}
ド・モアブルの定理（表\ref{bk:6-10:tab-kyoku}のべきの行で $r=1$）で $n=3$ とすると、
\[ (\cos\theta+j\sin\theta)^3=\cos3\theta+j\sin3\theta \]
左を、$(a+b)^3=a^3+3a^2b+3ab^2+b^3$（3-8）で展開します。$j^2=-1$、$j^3=-j$ を使うと、
\[ \cos^3\theta+3j\cos^2\theta\sin\theta-3\cos\theta\sin^2\theta-j\sin^3\theta \]
実部どうしを等しいとおき、$\sin^2\theta=1-\cos^2\theta$（4-8）を使うと、
\[ \cos3\theta=\cos^3\theta-3\cos\theta(1-\cos^2\theta)=4\cos^3\theta-3\cos\theta \]
$\theta=60°$ で確かめると、左は $\cos180°=-1$、右は $4\times\dfrac{1}{8}-3\times\dfrac{1}{2}=-1$。合っています。覚えにくい公式も、オイラーの公式から作れます。
\end{motto}

\Needspace{8\baselineskip}
\bksubsection{9}{交流は回る矢印：フェーザ}

\label{bk:6-10:phasor}家のコンセントの電圧は、$V_0\sin(\omega t+\phi)$ の形の交流です（5-15。$V_0$ は振幅〔V〕、$\phi$ は初期位相）。これを、回る矢印の影として見ます。
\begin{enumerate}[label={\bkcircled[\bkTF]{\arabic*}}, leftmargin=2.4\zw, itemsep=2pt, topsep=2pt]
\item 長さ $V_0$、角 $\phi$ の矢印 $V_0e^{j\phi}$ を用意する。
\item この矢印を、角速度 $\omega$ で回す。$t$ 秒後には $V_0e^{j\phi}\times e^{j\omega t}=V_0e^{j(\omega t+\phi)}$（肩のたし算）。
\item そのたての影（虚部）が、$V_0\sin(\omega t+\phi)$。
\end{enumerate}
①の矢印 $V_0e^{j\phi}$ を、交流の\textbf{フェーザ}といいます。\bkidx{ふぇーざ}{フェーザ}長さが振幅、角が初期位相を表します。「回り始める前の矢印」なので、時間 $t$ が入っていません。

\bkvikhead{同じ速さで回る矢印は、たし算できる}
同じ $\omega$ の交流が2つあるとき、2本の矢印は同じ速さで回るので、2本のあいだの角は変わりません。2本をたした矢印も、形を変えずに同じ速さで回ります。影は、矢印のたし算と同じようにたし算になる（たての成分どうしをたすだけ）ので、
\bkvikframe{同じ $\omega$ の波のたし算は、フェーザ（矢印）のたし算で計算できる}

\bkvikhead{4-9 のやってみようを、矢印で解く}
4-9 では、$\sin\theta+\sin(\theta+60°)$ を加法定理と合成で計算しました。矢印でやってみます。
\begin{enumerate}[label={\bkcircled[\bkTF]{\arabic*}}, leftmargin=2.4\zw, itemsep=2pt, topsep=2pt]
\item $\sin\theta$ のフェーザは、長さ1、角0の矢印：$1$。
\item $\sin(\theta+60°)$ のフェーザは、長さ1、角60°の矢印：$e^{j\pi/3}=0.5+0.866j$。
\item たす：$1+(0.5+0.866j)=1.5+0.866j$。
\item 大きさは $\sqrt{1.5^2+0.866^2}=\sqrt{2.25+0.75}=\sqrt3\approx1.73$。
\item 角は $\tan^{-1}\dfrac{0.866}{1.5}=\tan^{-1}0.577=30°$（右上の点）。
\end{enumerate}
だから $\sin\theta+\sin(\theta+60°)=\sqrt3\,\sin(\theta+30°)$。4-9 と同じ答えが、加法定理を使わずに、矢印のたし算だけで出ました。

\begin{bkfigure}
\begin{tikzpicture}[x=1mm, y=1mm]
  \begin{scope}
  \bkvikaxes{-4}{42}{-4}{26}
  \draw[bkBlue, line width=1.4pt, -{Stealth[length=4.5pt]}] (0,0) -- (20,0);
  \draw[bkOrange, line width=1.4pt, -{Stealth[length=4.5pt]}] (0,0) -- (60:20);
  \draw[bkOrange, dashed, line width=0.8pt, -{Stealth[length=4pt]}] (20,0) -- ++(60:20);
  \draw[bkPurple, line width=1.8pt, -{Stealth[length=5pt]}] (0,0) -- (30,17.32);
  \node[font=\sffamily\footnotesize, text=bkBlue, anchor=north, inner sep=1pt] at (10,-1) {1};
  \node[font=\sffamily\footnotesize, text=bkOrange!80!black, anchor=east, inner sep=1pt] at ($(60:20)+(-0.6,0)$) {$e^{j\pi/3}$};
  \node[font=\sffamily\footnotesize, text=bkPurple, anchor=west, inner sep=1pt, align=left] at (31,17.3) {和\\$\sqrt3e^{j\pi/6}$};
  \draw[bkInk, line width=0.5pt] (9,0) arc (0:30:9);
  \node[font=\sffamily\footnotesize, anchor=west, inner sep=1pt] at (11,2.6) {30°};
  \end{scope}
\end{tikzpicture}\hspace{5mm}%
\begin{tikzpicture}[baseline={(0,-6mm)}]
\begin{axis}[bkplot, scale only axis, width=62mm, height=32mm, xmin=0, xmax=375, ymin=-1.9, ymax=2.1, xtick={180,360}, xticklabels={180°,360°}, ytick=\empty,
    axis x line=bottom, x axis line style={draw=none}, xticklabel style={yshift=-0.5mm}, grid=none]
  \draw[bkInk, line width=0.8pt, -{Stealth[length=5pt]}] (axis cs:0,0) -- (axis cs:385,0) node[anchor=west, font=\sffamily\footnotesize, text=bkInk] {$\theta$};
  \draw[bkInk, line width=0.6pt] (axis cs:0,-1.9) -- (axis cs:0,2);
  \addplot[bkBlue, line width=1.2pt, domain=0:360, samples=90] {sin(x)};
  \addplot[bkOrange, line width=1.2pt, domain=0:360, samples=90] {sin(x+60)};
  \addplot[bkPurple, line width=1.7pt, domain=0:360, samples=90] {1.732*sin(x+30)};
  \node[font=\sffamily\footnotesize, text=bkPurple, anchor=south west, inner sep=1pt, fill=white] at (axis cs:75,1.75) {$\sqrt3\sin(\theta+30°)$};
\end{axis}
\end{tikzpicture}
\bkfigcaption{左：フェーザのたし算。オレンジの矢印を青の矢印の先へ平行に移す（点線）と、和（むらさき）の先に届く。右：青 $\sin\theta$ とオレンジ $\sin(\theta+60°)$ の和が、むらさきの波 $\sqrt3\sin(\theta+30°)$。}
\end{bkfigure}

\begin{seikaku}
ここでは、矢印のたての影（虚部）で $\sin$ を表しました。電気の本や、6-13 では、横の影（実部）で $\cos$ を表す約束を使うことが多いです。どちらの影を使っても、「同じ $\omega$ の波のたし算は、矢印のたし算」というきまりは同じです。ただし、1つの計算の中では、どちらか一方に決めて使います。
\end{seikaku}

\bkvikhead{三相交流：3本の矢印の和は0}
4-9 で、山の位置が120°ずつずれた3つの交流（三相交流）の和は、いつも0になることを、加法定理で確かめました。\bkidx{さんそうこうりゅう}{三相交流}フェーザなら、長さ1で、角が0、$-120°$、$-240°$ の3本の矢印です。
\begin{bkhenkei}[0.64]
1+e^{-j2\pi/3}+e^{-j4\pi/3} & 3本のフェーザのたし算 \\
=1+\Bigl(-\frac{1}{2}-\frac{\sqrt3}{2}j\Bigr)+\Bigl(-\frac{1}{2}+\frac{\sqrt3}{2}j\Bigr) & オイラーの公式（$-240°$ は $120°$ と同じ位置） \\
=\Bigl(1-\frac{1}{2}-\frac{1}{2}\Bigr)+\Bigl(-\frac{\sqrt3}{2}+\frac{\sqrt3}{2}\Bigr)j=0 & 実部どうし、虚部どうしをたす \\
\end{bkhenkei}
矢印を先から先へつないでいくと、正三角形ができて、もとの点にもどります。和の矢印が0なので、その影（3つの電圧の和）も、いつでも0です。
\begin{bkfigure}
\begin{tikzpicture}[x=1mm, y=1mm]
  \begin{scope}
  \bkvikaxes{-22}{27}{-21}{21}
  \draw[bkBlue, line width=1.4pt, -{Stealth[length=4.5pt]}] (0,0) -- (0:17);
  \draw[bkOrange, line width=1.4pt, -{Stealth[length=4.5pt]}] (0,0) -- (-120:17);
  \draw[bkPurple, line width=1.4pt, -{Stealth[length=4.5pt]}] (0,0) -- (120:17);
  \node[font=\sffamily\footnotesize, text=bkBlue, anchor=south, inner sep=1pt] at (12,0.8) {1};
  \node[font=\sffamily\footnotesize, text=bkOrange!80!black, anchor=north east, inner sep=1pt] at ($(-120:17)+(-0.3,0)$) {$e^{-j2\pi/3}$};
  \node[font=\sffamily\footnotesize, text=bkPurple, anchor=south east, inner sep=1pt] at ($(120:17)+(-0.3,0)$) {$e^{-j4\pi/3}$};
  \end{scope}
  \begin{scope}[xshift=62mm, yshift=6mm]
  \draw[bkBlue, line width=1.4pt, -{Stealth[length=4.5pt]}] (0,0) -- (17,0);
  \draw[bkOrange, line width=1.4pt, -{Stealth[length=4.5pt]}] (17,0) -- ++(-120:17);
  \draw[bkPurple, line width=1.4pt, -{Stealth[length=4.5pt]}] ($(17,0)+(-120:17)$) -- ++(120:17);
  \fill[bkInk] (0,0) circle (0.9);
  \node[font=\sffamily\footnotesize, anchor=west, inner sep=1pt, align=left] at (22,-7) {先から先へつなぐと\\もとの点にもどる\\（和は0）};
  \end{scope}
\end{tikzpicture}
\bkfigcaption{三相交流のフェーザ。左：120°ずつずれた3本の矢印。右：つなぐと正三角形になり、和は0。}
\end{bkfigure}

\Needspace{8\baselineskip}
\bksubsection{10}{微分は「$j\omega$ をかける」こと}

\label{bk:6-10:jomega}$e^{j\omega t}$ を微分します。5-14 のきまり $(e^{at})'=ae^{at}$ で $a=j\omega$ とすると（③で、$e^{jt}$ でも成り立つことを確かめた）、
\bkvikframe{$\dfrac{\dd}{\dd t}e^{j\omega t}=j\omega\,e^{j\omega t}$}
回る矢印を微分すると、$j\omega$ がかかります。$j\omega$ は「大きさ $\omega$、角90°」の数（$j\omega=\omega e^{j\pi/2}$）なので、
\bkvikline{\textbf{微分すると、矢印は $\omega$ 倍の長さになり、90°先へ進む}}
微分を、かけ算1回におきかえられるのです。
\begin{bkcenter}
\begin{bktable}{colspec={Q[l,m,wd=26mm] Q[c,m,wd=26mm] Q[l,m,wd=58mm]}, rowsep=2pt}
計算 & 矢印にかける数 & 意味 \\
微分 & $j\omega$ & $\omega$ 倍し、90°進める \\
2回微分 & $(j\omega)^2=-\omega^2$ & $\omega^2$ 倍し、180°（向きが反対） \\
積分 & $\dfrac{1}{j\omega}=-\dfrac{j}{\omega}$ & $\dfrac{1}{\omega}$ 倍し、90°おくらせる \\
\end{bktable}
\bktabcaption{2回微分の行は、5-15・5-18 の「2回微分すると、もとの $-\omega^2$ 倍」と同じ。}
\end{bkcenter}
積分の行の $\dfrac{1}{j}=-j$ は、分母と分子に $j$ をかけて、$\dfrac{j}{j^2}=\dfrac{j}{-1}=-j$（6-8）。$-j$ をかけると、右回りに90°回ります。

\bkvikhead{コンデンサーの電流は、電圧より90°進む}
コンデンサーにたまる電気の量 $q$〔C〕は、電圧 $v$〔V〕に比例します（$q=Cv$。$C$〔F〕は静電容量）。電流 $i$〔A〕は、$q$ の変わる速さです（6-4）。だから
\bkvikline{$i=\dfrac{\dd q}{\dd t}=C\dfrac{\dd v}{\dd t}$}
電圧が $v=V_0\sin\omega t$ なら、そのフェーザは $V_0$（角0の矢印）。微分は $j\omega$ をかけることなので、電流のフェーザは
\bkvikline{$j\omega C\times V_0=\omega CV_0\,e^{j\pi/2}$}
長さ $\omega CV_0$、角90°の矢印です。たての影をとると、
\bkvikline{$i=\omega CV_0\sin(\omega t+90°)=\omega CV_0\cos\omega t$}
電流は電圧より90°、つまり $\frac{1}{4}$ 周期先に進みます。5-15 で「位相のずれ（6-10）」と予告したことが、矢印1本で説明できました。

\begin{bkfigure}
\begin{tikzpicture}[x=1mm, y=1mm]
  \begin{scope}
  \bkvikaxes{-4}{30}{-4}{27}
  \draw[bkBlue, line width=1.5pt, -{Stealth[length=5pt]}] (0,0) -- (22,0);
  \draw[bkOrange, line width=1.5pt, -{Stealth[length=5pt]}] (0,0) -- (0,17);
  \draw[bkInk, line width=0.5pt] (6,0) arc (0:90:6);
  \node[font=\sffamily\footnotesize, anchor=south west, inner sep=1pt] at (4.6,4.6) {90°};
  \node[font=\sffamily\footnotesize, text=bkBlue, anchor=north, inner sep=1pt] at (13,-1) {電圧 $V_0$};
  \node[font=\sffamily\footnotesize, text=bkOrange!80!black, anchor=west, inner sep=1pt, align=left] at (1.8,13.5) {電流 $j\omega CV_0$};
  \end{scope}
  \begin{scope}[xshift=44mm, yshift=12mm]
  \draw[bkGray, line width=0.6pt, -{Stealth[length=3.5pt]}] (0,0) -- (72,0) node[anchor=west, font=\sffamily\footnotesize, text=bkInk] {$\omega t$};
  \draw[bkGray, line width=0.6pt] (0,-13) -- (0,13);
  \draw[bkBlue, line width=1.4pt, domain=0:6.5, samples=90] plot ({10*\x}, {10*sin(\x r)});
  \draw[bkOrange, line width=1.4pt, domain=0:6.5, samples=90] plot ({10*\x}, {7.5*cos(\x r)});
  \node[font=\sffamily\footnotesize, text=bkBlue, anchor=south, inner sep=1pt, fill=white] at (15.7,10.4) {電圧 $v$};
  \node[font=\sffamily\footnotesize, text=bkOrange!80!black, anchor=west, inner sep=1pt, fill=white] at (64,8.2) {電流 $i$};
  \draw[bkInk, line width=0.6pt, {Stealth[length=3pt]}-{Stealth[length=3pt]}] (0,-15) -- (15.7,-15);
  \node[font=\sffamily\footnotesize, anchor=north, inner sep=1pt] at (7.85,-15.8) {$\frac{1}{4}$ 周期};
  \end{scope}
\end{tikzpicture}
\bkfigcaption{コンデンサーの電圧と電流。左：フェーザ。電流の矢印は、電圧の矢印より90°先にある。右：電流の山（オレンジ）は、電圧の山（青）より $\frac{1}{4}$ 周期早く来る。}
\end{bkfigure}

\bkvikhead{コイルでは、電流が電圧より90°おくれる}
コイル（インダクタンス $L$〔H〕）の電圧は、電流の変わる速さに比例します（$v=L\dfrac{\dd i}{\dd t}$。6-4）。電流が $I_0\sin\omega t$ なら、電圧のフェーザは $j\omega L\times I_0$。電圧が電流より90°進む、言いかえると、電流が電圧より90°おくれます。\par
抵抗の $v=Ri$ と同じように、「電圧＝（かける数）×電流」の形にそろえてみます。この「かける数」を\textbf{インピーダンス}といいます。
\begin{itemize}[leftmargin=1.2\zw, itemsep=2pt, topsep=2pt, label={\color{\bkTT}●}]
\item コイル：電圧のフェーザ＝$j\omega L$×電流のフェーザ。かける数は $j\omega L$。
\item コンデンサー：上の式は「電流＝$j\omega C$×電圧」でした。両辺を $j\omega C$ でわって「電圧＝$\dfrac{1}{j\omega C}$×電流」。かける数は $\dfrac{1}{j\omega C}$。
\end{itemize}
コンデンサーの $j\omega C$ は「電圧から電流を求める数」、$\dfrac{1}{j\omega C}$ は「電流から電圧を求める数」で、向きが逆です。電圧と電流のどちらから求めるかを、いつも決めてから使いましょう。インピーダンス $j\omega L$・$\dfrac{1}{j\omega C}$ を使って共振を調べるのが 6-13 です。

\bksubsection{11}{肩に実部もあると：$e^{(\sigma+j\omega)t}$ はらせん}

\label{bk:6-10:rasen}肩が $(\sigma+j\omega)t$（$\sigma$ はシグマと読む実数）のときは、肩のたし算を逆に使って2つに分けます。
\bkvikline{$e^{(\sigma+j\omega)t}=e^{\sigma t}\times e^{j\omega t}$}
$e^{j\omega t}$ は長さ1で回る矢印、$e^{\sigma t}$ は実数で、矢印の長さを決めます。だから、この点は\textbf{回りながら、長さが $e^{\sigma t}$ で変わる}、つまり、らせんをえがきます。
\begin{itemize}[leftmargin=1.2\zw, itemsep=2pt, topsep=2pt, label={\color{\bkTT}●}]
\item $\sigma<0$：長さがへっていく（5-14 の $e^{-t/\tau}$）。内向きのらせん。
\item $\sigma=0$：長さ1のまま。円。
\item $\sigma>0$：長さがふえていく。外向きのらせん。
\end{itemize}
横の影（実部）は $e^{\sigma t}\cos\omega t$ です。$\sigma<0$ なら、「ゆれながら、だんだん小さくなる」波になります。5-18 で「ゆれながら止まる動きは、$e^{at}$ と $\sin$・$\cos$ のかけ算」と書いたのは、らせんの影だったからです。

\begin{bkfigure}
\begin{tikzpicture}[x=1mm, y=1mm]
  \begin{scope}
  \bkvikaxes{-25}{28}{-24}{25}
  \draw[bkBlue, line width=1.3pt, domain=0:12.57, samples=200] plot ({21*exp(-0.15*\x)*cos(\x r)}, {21*exp(-0.15*\x)*sin(\x r)});
  \fill[bkBlue] (21,0) circle (1);
  \node[font=\sffamily\footnotesize, anchor=north east, inner sep=1pt, text=bkBlue] at (20.4,-1.2) {$t=0$};
  \end{scope}
\end{tikzpicture}\hspace{6mm}%
\begin{tikzpicture}[baseline={(0,-24mm)}]
\begin{axis}[bkplot, scale only axis, width=64mm, height=40mm, xmin=0, xmax=4.2, ymin=-1.1, ymax=1.15, xtick={1,2,3,4}, ytick={-1,1},
    axis x line=bottom, x axis line style={draw=none}, xticklabel style={yshift=-0.5mm},
    yticklabel style={fill=white, inner sep=1pt, outer sep=2.3pt}]
  \draw[bkInk, line width=0.8pt, -{Stealth[length=5pt]}] (axis cs:0,0) -- (axis cs:4.3,0) node[anchor=west, font=\sffamily\footnotesize, text=bkInk] {$t$};
  \addplot[bkGray, dashed, line width=0.8pt, domain=0:4, samples=40] {exp(-x)};
  \addplot[bkGray, dashed, line width=0.8pt, domain=0:4, samples=40] {-exp(-x)};
  \addplot[bkPurple, line width=1.6pt, domain=0:4, samples=200] {exp(-x)*cos(deg(2*x))};
  \node[font=\sffamily\footnotesize, text=bkGray, anchor=south west, inner sep=1pt, fill=white] at (axis cs:1.05,0.40) {$e^{-t}$};
  \node[font=\sffamily\footnotesize, text=bkPurple, anchor=west, inner sep=1pt, fill=white] at (axis cs:1.75,-0.55) {$e^{-t}\cos2t$};
\end{axis}
\end{tikzpicture}
\bkfigcaption{左：$\sigma<0$ のときの $e^{(\sigma+j\omega)t}$ は、内向きのらせん。右：$e^{(-1+2j)t}$ の実部 $e^{-t}\cos2t$（むらさき）。山の高さは点線 $e^{-t}$ にそってへる。}\label{bk:6-10:fig-rasen}
\end{bkfigure}

\bkvikhead{3-10 の車の式とのつながり}
3-10 では、車のばねとダンパーの式から、特性方程式 $s^2+2s+5=0$ が出てきました。6-8 で解くと $s=-1\pm2j$ でした。$s=-1+2j$ なら、$e^{st}=e^{(-1+2j)t}$ の実部は $e^{-t}\cos2t$ で、図\ref{bk:6-10:fig-rasen}のむらさきのグラフです。「判別式が負なら、ゆれながらもどる」（3-10）のは、解が $\sigma+j\omega$ の形の複素数になり、$e^{st}$ がらせんになるからです。くわしくは 6-12 で解きます。

\begin{senmonsho}[電気回路の教科書ふうの書き方]
\begin{bkhonbun}
正弦波交流 $v(t)=\sqrt2\,V\sin(\omega t+\phi)$ を、フェーザ $\dot V=Ve^{j\phi}$ で表す。$v(t)$ の微分は $\dot V$ に $j\omega$ を乗ずることに、積分は $j\omega$ で除することに対応する。
\end{bkhonbun}
\bkyaku{$\sqrt2\,V$ が振幅（山の高さ）。$V$ は実効値（5-15。振幅を $\sqrt2$ でわった値）。日本の電気の本では、フェーザの大きさに実効値を使い、$\dot V$（ブイ・ドット）と書くことが多い。「乗ずる」はかける、「除する」はわる。電気の本では、$\dot V=V\angle\phi$ と書くこともある（6-9）。}
\end{senmonsho}





\begin{kotoba}
\small
\bkword{オイラーの公式}{$e^{j\theta}=\cos\theta+j\sin\theta$。指数関数と三角関数をつなぐ式。}
\bkword{オイラーの等式}{$e^{j\pi}+1=0$。オイラーの公式で $\theta=\pi$ としたもの。}
\bkword{フェーザ}{交流を表す、回り始める前の矢印（複素数）。長さが振幅、角が初期位相。}
\bkword{位相のずれ}{2つの波の山の位置のずれを、角で表したもの。フェーザでは、2本の矢印のあいだの角。}
\end{kotoba}



\begin{taisetsu}
\def\tsl#1{\makebox[7.4\zw][l]{\sffamily\bfseries\small\color{\bkTT}#1}}
\tsl{オイラーの公式}$e^{j\theta}=\cos\theta+j\sin\theta$（単位円の上の角 $\theta$ の点）\\[3pt]
\tsl{肩のたし算}$e^{ja}e^{jb}=e^{j(a+b)}$。$re^{j\theta}$ は大きさをかけ、角をたす\\[3pt]
\tsl{cos と sin}$\cos\theta=\dfrac{e^{j\theta}+e^{-j\theta}}{2}$\qquad $\sin\theta=\dfrac{e^{j\theta}-e^{-j\theta}}{2j}$\\[3pt]
\tsl{交流の合成}同じ $\omega$ の波のたし算は、フェーザのたし算\\[3pt]
\tsl{微分}$e^{j\omega t}$ の微分は、$j\omega$ をかける（$\omega$ 倍して90°進める）
\end{taisetsu}







\bkkigocards{\bkkigocard{$e^{j\theta}$}{イー・ジェイ・シータ（$e$ の $j\theta$ 乗）}{単位円の上の角 $\theta$ の点。$\cos\theta+j\sin\theta$}{complex exponential}{6-10}}
{\bkkigocard{$j\omega$}{ジェイ・オメガ}{$e^{j\omega t}$ を微分すると前に出る数。かけると $\omega$ 倍し90°進める}{j-omega}{6-10}}



\begin{column}{オイラーという人}
レオンハルト・オイラー（1707〜1783年）は、スイスのバーゼルに生まれ、ロシアのサンクトペテルブルクとドイツのベルリンで研究した数学者です。$e$、$\pi$、$f(x)$、$\sum$ などの記号を広めたのもオイラーです。後の人生では、ほとんど目が見えなくなりましたが、計算を頭の中で進め、人に書きとってもらって、研究を続けました。一生のあいだに書いた論文や本は、800をこえます。6-2 のオイラー法も、この章のオイラーの公式も、同じオイラーの仕事です。
\end{column}

% Ver2（No.34）：例題の前のページの下の空きを、第7巻（7-6 回転行列）へのつながりで埋めた
\begin{motto}[4]
\bkvikbox{$e^{j\theta}$ をかける回転を、ふつうの数の式で書く}
点 $(x,\ y)$ を $x+yj$ と考えて、$e^{j\theta}=\cos\theta+j\sin\theta$ をかけます。$j^2=-1$ を使って広げると、
\[ (\cos\theta+j\sin\theta)(x+yj)=(x\cos\theta-y\sin\theta)+(x\sin\theta+y\cos\theta)\,j \]
です。だから、点 $(x,\ y)$ を $\theta$ 回した点は、$(x\cos\theta-y\sin\theta,\ \ x\sin\theta+y\cos\theta)$ です。たとえば $\theta=90^\circ$ なら $\cos\theta=0$、$\sin\theta=1$ で、$(x,\ y)$ は $(-y,\ x)$ に移ります（6-8 の「$j$ をかけると90°回る」と同じ）。

第7巻（7-6）では、この2つの式を「回転行列」という表の形のかけ算で書きます。「30°回して45°回すと、75°回したのと同じ」は、ここでは $e^{j30^\circ}e^{j45^\circ}=e^{j75^\circ}$（指数法則）として、もう見えています。
\end{motto}

%% ---------------- 例題 ----------------
\bksection{reidai}{例題でたしかめよう}

\begin{reidai}[$e^{j\theta}$ を $a+bj$ に直す]
次の数を $a+bj$ の形に直しましょう。\par（1）$e^{j\pi/3}$\quad（2）$e^{j3\pi/4}$\quad（3）$4e^{-j\pi/2}$\par
\kangaekata オイラーの公式 $e^{j\theta}=\cos\theta+j\sin\theta$ に入れる。角は単位円の上で読む（ラジアンを度に直すと分かりやすい：$\pi$ rad が180°）。マイナスの角は右回り。\par
\shiki\hangindent=4.3\zw\hangafter=1 （1）$\dfrac{\pi}{3}$ は60°。$\cos60°+j\sin60°=\dfrac{1}{2}+\dfrac{\sqrt3}{2}j$\par
\bkvikq （2）$\dfrac{3\pi}{4}$ は135°（左上）。横は負、たては正。$\cos135°=-\dfrac{\sqrt2}{2}$、$\sin135°=\dfrac{\sqrt2}{2}$\par
\bkvikq （3）$-\dfrac{\pi}{2}$ は、右回りに90°（真下）。\par
\bkvikb $e^{-j\pi/2}=\cos(-90°)+j\sin(-90°)=0-j=-j$\par
\bkvikb 4倍して $4\times(-j)=-4j$\par
\par\vspace{3pt}{\centering
\begin{tikzpicture}[x=1mm, y=1mm]
  \draw[bkGray, line width=0.6pt, -{Stealth[length=3pt]}] (-20,0) -- (24,0) node[anchor=north east, font=\sffamily\footnotesize, inner sep=1pt, yshift=-1mm] {実部};
  \draw[bkGray, line width=0.6pt, -{Stealth[length=3pt]}] (0,-16) -- (0,18) node[anchor=north west, font=\sffamily\footnotesize, inner sep=1pt, xshift=1mm] {虚部};
  \draw[bkInk!40, line width=0.7pt] (0,0) circle (13);
  \foreach \a/\c in {60/bkBlue, 135/bkOrange, -90/bkPurple} {\draw[\c, line width=1.2pt, -{Stealth[length=4pt]}] (0,0) -- (\a:13);}
  \node[font=\sffamily\footnotesize, inner sep=1pt, text=bkBlue, anchor=west] at (8,10) {（1）60°};
  \node[font=\sffamily\footnotesize, inner sep=1pt, text=bkOrange!80!black, anchor=south east] at (135:13.5) {（2）135°};
  \node[font=\sffamily\footnotesize, inner sep=1pt, text=bkPurple, anchor=west] at (1.5,-11) {（3）$-90°$（長さを4倍する）};
\end{tikzpicture}\par}\vspace{2pt}
\kotae\hangindent=4.3\zw\hangafter=1 （1）$\dfrac{1}{2}+\dfrac{\sqrt3}{2}j$\quad（2）$-\dfrac{\sqrt2}{2}+\dfrac{\sqrt2}{2}j$\quad（3）$-4j$
\end{reidai}

\begin{jibun}
次の数を $a+bj$ の形に直しましょう。\par（1）$e^{j\pi/6}$\quad（2）$e^{j2\pi/3}$\quad（3）$2e^{-j\pi}$
\bkjans{（1）\nobreak $\dfrac{\sqrt3}{2}+\dfrac{1}{2}j$。（2）\nobreak 120°なので $-\dfrac{1}{2}+\dfrac{\sqrt3}{2}j$。（3）\nobreak $e^{-j\pi}=-1$ なので $-2$。}
\end{jibun}

\begin{reidai}[$a+bj$ を $re^{j\theta}$ に直す]
次の数を $re^{j\theta}$ の形に直しましょう（$-\pi<\theta\leqq\pi$）。\par（1）$-1+\sqrt3\,j$\quad（2）$1-j$\par
\kangaekata $r$ は原点からの距離（三平方）。$\theta$ は、点がどの区画にあるかを図でたしかめてから決める（6-9）。\par
\shiki\hangindent=4.3\zw\hangafter=1 （1）$r=\sqrt{(-1)^2+(\sqrt3)^2}=\sqrt{1+3}=2$\par
\bkvikb $\cos\theta=\dfrac{-1}{2}$、$\sin\theta=\dfrac{\sqrt3}{2}$。左上の点なので $\theta=120°=\dfrac{2\pi}{3}$\par
\bkvikq （2）$r=\sqrt{1^2+(-1)^2}=\sqrt2$\par
\bkvikb $\cos\theta=\dfrac{1}{\sqrt2}$、$\sin\theta=-\dfrac{1}{\sqrt2}$。右下の点なので $\theta=-45°=-\dfrac{\pi}{4}$\par
\par\vspace{3pt}{\centering
\begin{tikzpicture}[x=1mm, y=1mm]
  \draw[bkGray, line width=0.6pt, -{Stealth[length=3pt]}] (-18,0) -- (27,0) node[anchor=north east, font=\sffamily\footnotesize, inner sep=1pt, yshift=-1mm] {実部};
  \draw[bkGray, line width=0.6pt, -{Stealth[length=3pt]}] (0,-16) -- (0,22) node[anchor=north west, font=\sffamily\footnotesize, inner sep=1pt, xshift=1mm] {虚部};
  \draw[bkBlue, line width=1.2pt, -{Stealth[length=4pt]}] (0,0) -- (-8,13.86);
  \draw[bkOrange, line width=1.2pt, -{Stealth[length=4pt]}] (0,0) -- (8,-8);
  \draw[bkBlue, line width=0.5pt] (4,0) arc (0:120:4);
  \draw[bkOrange, line width=0.5pt] (6,0) arc (0:-45:6);
  \node[font=\sffamily\footnotesize, inner sep=1pt, text=bkBlue, anchor=south east] at (-8.3,14) {（1）$-1+\sqrt3\,j$};
  \node[font=\sffamily\footnotesize, inner sep=1pt, text=bkBlue, anchor=south west] at (2.5,3.6) {120°};
  \node[font=\sffamily\footnotesize, inner sep=1pt, text=bkOrange!80!black, anchor=west] at (8.6,-8.4) {（2）$1-j$};
  \node[font=\sffamily\footnotesize, inner sep=1pt, text=bkOrange!80!black, anchor=west] at (7.5,-4.6) {$-45°$};
\end{tikzpicture}\par}\vspace{2pt}
\kotae\hangindent=4.3\zw\hangafter=1 （1）$2e^{j2\pi/3}$\quad（2）$\sqrt2\,e^{-j\pi/4}$
\end{reidai}

\begin{jibun}
次の数を $re^{j\theta}$ の形に直しましょう（$-\pi<\theta\leqq\pi$）。\par（1）$\sqrt3-j$\quad（2）$-2+2j$
\bkjans{（1）\nobreak $r=2$。右下で $\theta=-30°$。$2e^{-j\pi/6}$。（2）\nobreak $r=\sqrt8=2\sqrt2$。左上で $\theta=135°$。$2\sqrt2\,e^{j3\pi/4}$。}
\end{jibun}

\begin{reidai}[極形式のかけ算・わり算]
次の計算をし、答えを $re^{j\theta}$ と $a+bj$ の両方で書きましょう。\par
（1）$2e^{j\pi/6}\times3e^{j\pi/3}$\quad（2）$\dfrac{6e^{j\pi/2}}{2e^{j\pi/6}}$\par
\kangaekata かけ算は大きさをかけて角（肩）をたす。わり算は大きさをわって角をひく（⑦の表）。\par
\shiki\hangindent=4.3\zw\hangafter=1 （1）大きさ $2\times3=6$。角 $\dfrac{\pi}{6}+\dfrac{\pi}{3}=\dfrac{\pi}{6}+\dfrac{2\pi}{6}=\dfrac{3\pi}{6}=\dfrac{\pi}{2}$\par
\bkvikb $6e^{j\pi/2}=6(\cos90°+j\sin90°)=6j$\par
\bkvikq （2）大きさ $6\div2=3$。角 $\dfrac{\pi}{2}-\dfrac{\pi}{6}=\dfrac{3\pi}{6}-\dfrac{\pi}{6}=\dfrac{2\pi}{6}=\dfrac{\pi}{3}$\par
\bkvikb $3e^{j\pi/3}=3\Bigl(\dfrac{1}{2}+\dfrac{\sqrt3}{2}j\Bigr)=\dfrac{3}{2}+\dfrac{3\sqrt3}{2}j$\par
\par\vspace{3pt}{\centering
\begin{tikzpicture}[x=1mm, y=1mm]
  \draw[bkGray, line width=0.6pt, -{Stealth[length=3pt]}] (-4,0) -- (30,0) node[anchor=north east, font=\sffamily\footnotesize, inner sep=1pt, yshift=-1mm] {実部};
  \draw[bkGray, line width=0.6pt, -{Stealth[length=3pt]}] (0,-4) -- (0,28) node[anchor=north west, font=\sffamily\footnotesize, inner sep=1pt, xshift=1mm] {虚部};
  \draw[bkBlue, line width=1.2pt, -{Stealth[length=4pt]}] (0,0) -- (30:8);
  \draw[bkOrange, line width=1.2pt, -{Stealth[length=4pt]}] (0,0) -- (60:12);
  \draw[bkPurple, line width=1.6pt, -{Stealth[length=4.5pt]}] (0,0) -- (90:24);
  \node[font=\sffamily\footnotesize, inner sep=1pt, text=bkBlue, anchor=west] at ($(30:8)+(0.8,-0.4)$) {長さ2、30°};
  \node[font=\sffamily\footnotesize, inner sep=1pt, text=bkOrange!80!black, anchor=west] at ($(60:12)+(0.8,0)$) {長さ3、60°};
  \node[font=\sffamily\footnotesize, inner sep=1pt, text=bkPurple, anchor=west] at (1.5,21) {かけると 長さ6、90°};
\end{tikzpicture}\par}\vspace{2pt}
\kotae\hangindent=4.3\zw\hangafter=1 （1）$6e^{j\pi/2}=6j$\quad（2）$3e^{j\pi/3}=\dfrac{3}{2}+\dfrac{3\sqrt3}{2}j$
\end{reidai}

\begin{jibun}
次の計算をし、$re^{j\theta}$ と $a+bj$ の両方で書きましょう。\par
（1）$4e^{j\pi/12}\times2e^{j\pi/4}$\quad（2）$\dfrac{10e^{j5\pi/6}}{5e^{j\pi/2}}$
\bkjans{（1）\nobreak $8e^{j\pi/3}=4+4\sqrt3\,j$。（2）\nobreak $2e^{j\pi/3}=1+\sqrt3\,j$。}
\end{jibun}

\begin{reidai}[べきを計算する]
$(\sqrt3+j)^6$ を計算しましょう。\par
\kangaekata 6回かけて展開するのは大変。極形式に直してから、大きさを6乗、角を6倍する（⑦）。\par
\shiki\hangindent=2.6\zw\hangafter=1 $\sqrt3+j$ の大きさは $\sqrt{3+1}=2$。右上の点で $\cos\theta=\dfrac{\sqrt3}{2}$ なので角は $\dfrac{\pi}{6}$（30°）\par
\bkvika $(\sqrt3+j)^6=\bigl(2e^{j\pi/6}\bigr)^6=2^6\,e^{j\,6\times\pi/6}=64\,e^{j\pi}$\par
\bkvika $e^{j\pi}=-1$ なので $64\times(-1)$\par
\par\vspace{3pt}{\centering
\begin{tikzpicture}[x=1mm, y=1mm]
  \draw[bkGray, line width=0.6pt, -{Stealth[length=3pt]}] (-20,0) -- (22,0) node[anchor=north east, font=\sffamily\footnotesize, inner sep=1pt, yshift=-1mm] {実部};
  \draw[bkGray, line width=0.6pt, -{Stealth[length=3pt]}] (0,-3) -- (0,18) node[anchor=north west, font=\sffamily\footnotesize, inner sep=1pt, xshift=1mm] {虚部};
  \draw[bkInk!40, line width=0.7pt] (13,0) arc (0:180:13);
  \foreach \k in {0,...,6} {\fill[bkBlue] ({30*\k}:13) circle (0.9); \node[font=\sffamily\footnotesize, inner sep=1pt, anchor={180+30*\k}] at ({30*\k}:14.5) {\k};}
  \node[font=\sffamily\footnotesize, inner sep=1pt, anchor=north, text=bkBlue] at (0,-1.5) {1回かけるごとに30°ずつ回る（数字はかけた回数）};
\end{tikzpicture}\par}\vspace{2pt}
\kotae\hangindent=2.6\zw\hangafter=1 $-64$
\end{reidai}

\begin{jibun}
$(1-j)^4$ を計算しましょう。
\bkjans{$1-j=\sqrt2\,e^{-j\pi/4}$。$(\sqrt2)^4e^{-j\pi}=4\times(-1)=-4$。}
\end{jibun}

\begin{reidai}[$\cos$・$\sin$ を $e^{j\theta}$ で計算する]
$\theta=\dfrac{\pi}{6}$ のとき、$\dfrac{e^{j\theta}+e^{-j\theta}}{2}$ と $\dfrac{e^{j\theta}-e^{-j\theta}}{2j}$ を計算し、$\cos\dfrac{\pi}{6}$、$\sin\dfrac{\pi}{6}$ と同じになることを確かめましょう。\par
\kangaekata $e^{j\pi/6}$ を $a+bj$ に直し、$e^{-j\pi/6}$ はその共役（虚部の符号を反対に）にする（④）。\par
\shiki\hangindent=2.6\zw\hangafter=1 $e^{j\pi/6}=\dfrac{\sqrt3}{2}+\dfrac{1}{2}j$、$e^{-j\pi/6}=\dfrac{\sqrt3}{2}-\dfrac{1}{2}j$\par
\bkvika たす：$\dfrac{\sqrt3}{2}+\dfrac{\sqrt3}{2}=\sqrt3$（虚部は打ち消しあう）。2でわって $\dfrac{\sqrt3}{2}$\par
\bkvika ひく：$\dfrac{1}{2}j-\Bigl(-\dfrac{1}{2}j\Bigr)=j$（実部は打ち消しあう）。$2j$ でわって $\dfrac{j}{2j}=\dfrac{1}{2}$\par
\par\vspace{3pt}{\centering
\begin{tikzpicture}[x=1mm, y=1mm]
  \draw[bkGray, line width=0.6pt, -{Stealth[length=3pt]}] (-4,0) -- (32,0) node[anchor=north east, font=\sffamily\footnotesize, inner sep=1pt, yshift=-1mm] {実部};
  \draw[bkGray, line width=0.6pt, -{Stealth[length=3pt]}] (0,-12) -- (0,13) node[anchor=north west, font=\sffamily\footnotesize, inner sep=1pt, xshift=1mm] {虚部};
  \draw[bkBlue, line width=1.2pt, -{Stealth[length=4pt]}] (0,0) -- (30:14);
  \draw[bkOrange, line width=1.2pt, -{Stealth[length=4pt]}] (0,0) -- (-30:14);
  \draw[bkGray, dashed, line width=0.5pt] (30:14) -- (24.25,0) -- (-30:14);
  \draw[bkPurple, line width=1.6pt, -{Stealth[length=4.5pt]}] (0,0) -- (24.25,0);
  \node[font=\sffamily\footnotesize, inner sep=1pt, text=bkBlue, anchor=south west] at (30:14.3) {$e^{j\pi/6}$};
  \node[font=\sffamily\footnotesize, inner sep=1pt, text=bkOrange!80!black, anchor=north west] at (-30:14.3) {$e^{-j\pi/6}$};
  \node[font=\sffamily\footnotesize, inner sep=1pt, text=bkPurple, anchor=south] at (17,0.8) {和 $\sqrt3$};
\end{tikzpicture}\par}\vspace{2pt}
\kotae\hangindent=2.6\zw\hangafter=1 $\dfrac{\sqrt3}{2}$ と $\dfrac{1}{2}$。$\cos\dfrac{\pi}{6}=\dfrac{\sqrt3}{2}$、$\sin\dfrac{\pi}{6}=\dfrac{1}{2}$ と同じ
\end{reidai}

\begin{jibun}
$\theta=\dfrac{\pi}{4}$ のとき、$\dfrac{e^{j\theta}+e^{-j\theta}}{2}$ と $\dfrac{e^{j\theta}-e^{-j\theta}}{2j}$ を計算しましょう。
\bkjans{$e^{\pm j\pi/4}=\dfrac{\sqrt2}{2}\pm\dfrac{\sqrt2}{2}j$。たして2でわると $\dfrac{\sqrt2}{2}$。ひくと $\sqrt2\,j$、$2j$ でわると $\dfrac{\sqrt2}{2}$。}
\end{jibun}

\begin{reidai}[交流の電圧をフェーザでたす]
2つの交流の電源を直列につなぎました。電圧は $v_1=100\sin\omega t$〔V〕と $v_2=100\sin(\omega t+60°)$〔V〕です。合計の電圧を、1つの $\sin$ の形で表しましょう。\par
\kangaekata 同じ $\omega$ なので、フェーザ（長さ＝振幅、角＝初期位相）をたし、和の矢印の長さと角を読む（⑨）。\par
\shiki\hangindent=2.6\zw\hangafter=1 $v_1$ のフェーザ：$100$。$v_2$ のフェーザ：$100e^{j\pi/3}=50+86.6j$\par
\bkvika 和：$100+50+86.6j=150+86.6j$\par
\bkvika 大きさ：$\sqrt{150^2+86.6^2}=\sqrt{22500+7500}=\sqrt{30000}\approx173$\par
\bkvika 角：$\tan^{-1}\dfrac{86.6}{150}=\tan^{-1}0.577=30°$（右上の点）\par
\par\vspace{3pt}{\centering
\begin{tikzpicture}[x=1mm, y=1mm]
  \draw[bkGray, line width=0.6pt, -{Stealth[length=3pt]}] (-4,0) -- (44,0) node[anchor=north east, font=\sffamily\footnotesize, inner sep=1pt, yshift=-1mm] {実部};
  \draw[bkGray, line width=0.6pt, -{Stealth[length=3pt]}] (0,-4) -- (0,22) node[anchor=north west, font=\sffamily\footnotesize, inner sep=1pt, xshift=1mm] {虚部};
  \draw[bkBlue, line width=1.2pt, -{Stealth[length=4pt]}] (0,0) -- (20,0);
  \draw[bkOrange, dashed, line width=0.9pt, -{Stealth[length=4pt]}] (20,0) -- ++(60:20);
  \draw[bkPurple, line width=1.6pt, -{Stealth[length=4.5pt]}] (0,0) -- (30,17.32);
  \node[font=\sffamily\footnotesize, inner sep=1pt, text=bkBlue, anchor=north] at (10,-1) {$v_1$：100};
  \node[font=\sffamily\footnotesize, inner sep=1pt, text=bkOrange!80!black, anchor=west] at (26,7) {$v_2$：100、60°};
  \node[font=\sffamily\footnotesize, inner sep=1pt, text=bkPurple, anchor=east] at (13,10) {和 173、30°};
\end{tikzpicture}\par}\vspace{2pt}
\kotae\hangindent=2.6\zw\hangafter=1 $v_1+v_2\approx173\sin(\omega t+30°)$〔V〕（山の位置が60°ずれているので、振幅は200ではない）
\end{reidai}

\begin{jibun}
$v_1=100\sin\omega t$〔V〕と $v_2=100\sin(\omega t+120°)$〔V〕の和を、1つの $\sin$ の形で表しましょう。
\bkjans{$100+100e^{j2\pi/3}=100-50+86.6j=50+86.6j$。大きさ $\sqrt{2500+7500}=100$、角60°。$100\sin(\omega t+60°)$〔V〕。}
\end{jibun}

\begin{reidai}[$\sin$ と $\cos$ の和をフェーザで]
$3\sin\omega t+4\cos\omega t$ を、1つの $\sin$ の形で表しましょう（4-9 の例題と同じ数）。\par
\kangaekata $\cos\omega t=\sin(\omega t+90°)$（4-8）なので、$4\cos\omega t$ のフェーザは長さ4、角90°の矢印 $4j$。\par
\shiki\hangindent=2.6\zw\hangafter=1 フェーザ：$3$ と $4e^{j\pi/2}=4j$。和は $3+4j$\par
\bkvika 大きさ：$\sqrt{3^2+4^2}=\sqrt{25}=5$\par
\bkvika 角：$\tan^{-1}\dfrac{4}{3}\approx53.1°$（右上の点）\par
\par\vspace{3pt}{\centering
\begin{tikzpicture}[x=1mm, y=1mm]
  \draw[bkGray, line width=0.6pt, -{Stealth[length=3pt]}] (-4,0) -- (30,0) node[anchor=north east, font=\sffamily\footnotesize, inner sep=1pt, yshift=-1mm] {実部};
  \draw[bkGray, line width=0.6pt, -{Stealth[length=3pt]}] (0,-4) -- (0,26) node[anchor=north west, font=\sffamily\footnotesize, inner sep=1pt, xshift=1mm] {虚部};
  \draw[bkBlue, line width=1.2pt, -{Stealth[length=4pt]}] (0,0) -- (12,0);
  \draw[bkOrange, dashed, line width=0.9pt, -{Stealth[length=4pt]}] (12,0) -- (12,16);
  \draw[bkPurple, line width=1.6pt, -{Stealth[length=4.5pt]}] (0,0) -- (12,16);
  \node[font=\sffamily\footnotesize, inner sep=1pt, text=bkBlue, anchor=north] at (6,-1) {3};
  \node[font=\sffamily\footnotesize, inner sep=1pt, text=bkOrange!80!black, anchor=west] at (12.8,8) {$4j$};
  \node[font=\sffamily\footnotesize, inner sep=1pt, text=bkPurple, anchor=east] at (5.4,10) {和 5};
  \draw[bkInk, line width=0.5pt] (5,0) arc (0:53.1:5);
  \node[font=\sffamily\footnotesize, inner sep=1pt, anchor=west] at (5.8,1.8) {53.1°};
\end{tikzpicture}\par}\vspace{2pt}
\kotae\hangindent=2.6\zw\hangafter=1 $5\sin(\omega t+53.1°)$
\end{reidai}

\begin{jibun}
$4\sin\omega t+3\cos\omega t$ を、1つの $\sin$ の形で表しましょう。
\bkjans{フェーザの和 $4+3j$。大きさ5、角 $\tan^{-1}\dfrac{3}{4}\approx36.9°$。$5\sin(\omega t+36.9°)$。}
\end{jibun}

\begin{reidai}[微分・積分を $j\omega$ で計算する]
$e^{j2t}$ を、1回微分、2回微分、積分（積分定数は考えない）しましょう。また、それぞれ、もとの矢印 $e^{j2t}$ を何倍し、何度回したものかを答えましょう。\par
\kangaekata $\omega=2$。微分は $j\omega=2j$ をかける。積分は $j\omega$ でわる（⑩の表）。\par
\shiki\hangindent=2.6\zw\hangafter=1 微分：$2j\,e^{j2t}$。$2j=2e^{j\pi/2}$ なので、2倍して90°進めたもの\par
\bkvika 2回微分：$(2j)^2e^{j2t}=4j^2e^{j2t}=-4e^{j2t}$。4倍して180°回したもの\par
\bkvika 積分：$\dfrac{e^{j2t}}{2j}=\dfrac{j}{2j^2}\,e^{j2t}=-\dfrac{j}{2}\,e^{j2t}$。$\dfrac{1}{2}$ 倍して90°おくらせたもの\par
\bkvika 確かめ：$\Bigl(-\dfrac{j}{2}e^{j2t}\Bigr)'=-\dfrac{j}{2}\times2j\,e^{j2t}=-j^2e^{j2t}=e^{j2t}$\par
\par\vspace{3pt}{\centering
\begin{tikzpicture}[x=2mm, y=2mm]
  \draw[bkGray, line width=0.6pt, -{Stealth[length=3pt]}] (-25,0) -- (14,0) node[anchor=north east, font=\sffamily\footnotesize, inner sep=1pt, yshift=-1mm] {実部};
  \draw[bkGray, line width=0.6pt, -{Stealth[length=3pt]}] (0,-5) -- (0,12) node[anchor=north west, font=\sffamily\footnotesize, inner sep=1pt, xshift=1mm] {虚部};
  \draw[bkBlue, line width=1.3pt, -{Stealth[length=4pt]}] (0,0) -- (5,0);
  \draw[bkOrange, line width=1.3pt, -{Stealth[length=4pt]}] (0,0) -- (0,10);
  \draw[bkPurple, line width=1.3pt, -{Stealth[length=4pt]}] (0,0) -- (-20,0);
  \draw[bkGreen!70!black, line width=1.3pt, -{Stealth[length=4pt]}] (0,0) -- (0,-2.5);
  \node[font=\sffamily\footnotesize, inner sep=1pt, text=bkBlue, anchor=south west] at (1,0.6) {もと（長さ1）};
  \node[font=\sffamily\footnotesize, inner sep=1pt, text=bkOrange!80!black, anchor=west] at (0.6,7) {微分（2倍、90°進む）};
  \node[font=\sffamily\footnotesize, inner sep=1pt, text=bkPurple, anchor=south] at (-17,0.6) {2回微分（4倍、180°）};
  \node[font=\sffamily\footnotesize, inner sep=1pt, text=bkGreen!50!black, anchor=east] at (-0.6,-3) {積分（$\frac{1}{2}$倍、90°おくれる）};
\end{tikzpicture}\par}\vspace{2pt}
\kotae\hangindent=2.6\zw\hangafter=1 微分 $2je^{j2t}$（2倍、90°進む）、2回微分 $-4e^{j2t}$（4倍、180°）、積分 $-\dfrac{j}{2}e^{j2t}$（$\dfrac{1}{2}$ 倍、90°おくれる）
\end{reidai}

\begin{jibun}
$e^{j5t}$ を、1回微分、2回微分、積分（積分定数は考えない）しましょう。それぞれ何倍し、何度回したものですか。
\bkjans{微分 $5je^{j5t}$（5倍、90°進む）。2回微分 $-25e^{j5t}$（25倍、180°）。積分 $\dfrac{e^{j5t}}{5j}=-\dfrac{j}{5}e^{j5t}$（$\dfrac{1}{5}$ 倍、90°おくれる）。}
\end{jibun}

\begin{reidai}[コンデンサーに流れる交流の電流]
静電容量 \qty{10}{\micro F} のコンデンサーに、$v=100\sin\omega t$〔V〕、周波数 \qty{50}{Hz} の電圧をかけます。電流の振幅と、電圧とのずれを求めましょう（$\pi=3.14$）。\par
\kangaekata 電流のフェーザは、電圧のフェーザに $j\omega C$ をかけたもの（⑩）。$\omega=2\pi f$（4-8）。\par
\shiki\hangindent=2.6\zw\hangafter=1 $\omega=2\pi\times50=2\times3.14\times50=314$〔rad/s〕\par
\bkvika $C=\qty{10}{\micro F}=10\times10^{-6}$\,F$=1\times10^{-5}$\,F\par
\bkvika 電流のフェーザ：$j\omega C\times100=j\times314\times10^{-5}\times100=0.314j$\par
\bkvika 長さ0.314、角90°の矢印\par
\par\vspace{3pt}{\centering
\begin{tikzpicture}[x=1mm, y=1mm]
  \draw[bkGray, line width=0.6pt, -{Stealth[length=3pt]}] (-4,0) -- (40,0) node[anchor=north east, font=\sffamily\footnotesize, inner sep=1pt, yshift=-1mm] {実部};
  \draw[bkGray, line width=0.6pt, -{Stealth[length=3pt]}] (0,-4) -- (0,22) node[anchor=north west, font=\sffamily\footnotesize, inner sep=1pt, xshift=1mm] {虚部};
  \draw[bkBlue, line width=1.4pt, -{Stealth[length=4.5pt]}] (0,0) -- (30,0);
  \draw[bkOrange, line width=1.4pt, -{Stealth[length=4.5pt]}] (0,0) -- (0,15);
  \node[font=\sffamily\footnotesize, inner sep=1pt, text=bkBlue, anchor=north] at (16,-1) {電圧 100\,V};
  \node[font=\sffamily\footnotesize, inner sep=1pt, text=bkOrange!80!black, anchor=west] at (1.2,12) {電流 0.314\,A（90°進む）};
\end{tikzpicture}\par}\vspace{2pt}
\kotae\hangindent=2.6\zw\hangafter=1 振幅 約\qty{0.31}{A}。電流は電圧より90°（$\frac{1}{4}$ 周期）進む
\end{reidai}

\begin{jibun}
静電容量 \qty{20}{\micro F} のコンデンサーに、$v=100\sin\omega t$〔V〕、周波数 \qty{60}{Hz} の電圧をかけます。電流の振幅と、電圧とのずれを求めましょう（$\pi=3.14$）。
\bkjans{$\omega=2\times3.14\times60\approx377$〔rad/s〕。$377\times20\times10^{-6}\times100\approx0.754$。振幅 約\qty{0.75}{A}、電流は電圧より90°進む。}
\end{jibun}

\begin{reidai}[コイルに流れる交流の電流：インピーダンスでわる]
インダクタンス \qty{0.1}{H} のコイルに、$v=100\sin\omega t$〔V〕、周波数 \qty{50}{Hz} の電圧をかけます。電流の振幅と、電圧とのずれを求めましょう（$\pi=3.14$）。\par
\kangaekata コイルは「電圧＝$j\omega L$×電流」（⑩のインピーダンス）。電流を求めるので、電圧のフェーザを $j\omega L$ で\textbf{わる}。\par
\shiki\hangindent=2.6\zw\hangafter=1 $\omega=2\pi\times50=314$〔rad/s〕。$j\omega L=j\times314\times0.1=31.4j$〔Ω〕\par
\bkvika 電流のフェーザ：$\dfrac{100}{31.4j}=\dfrac{100}{31.4}\times\dfrac{1}{j}=3.18\times(-j)=-3.18j$\par
\bkvika （$\dfrac{1}{j}=-j$：⑩の表の下）。長さ3.18、角 $-90°$ の矢印\par
\kotae\hangindent=2.6\zw\hangafter=1 振幅 約\qty{3.2}{A}。電流は電圧より90°（$\frac{1}{4}$ 周期）おくれる
\end{reidai}

\begin{jibun}
インダクタンス \qty{0.2}{H} のコイルに、$v=100\sin\omega t$〔V〕、周波数 \qty{60}{Hz} の電圧をかけます。電流の振幅と、電圧とのずれを求めましょう（$\pi=3.14$）。
\bkjans{$\omega=2\times3.14\times60\approx377$〔rad/s〕。$j\omega L=377\times0.2j\approx75.4j$〔Ω〕。$\dfrac{100}{75.4j}\approx-1.33j$。振幅 約\qty{1.3}{A}、電流は電圧より90°おくれる。}
\end{jibun}

\begin{reidai}[$e^{(\sigma+j\omega)t}$ の実部]
$e^{(-1+2j)t}$ の実部を、$t$ の式で書きましょう。また、$t=0$、$\dfrac{\pi}{4}$、$\dfrac{\pi}{2}$ のときの値を求めましょう（$e^{-\pi/2}=0.208$）。\par
\kangaekata 肩のたし算を逆に使って $e^{-t}\times e^{j2t}$ に分ける。$e^{-t}$ は実数なので、実部は $e^{-t}\times\cos2t$（⑪）。\par
\shiki\hangindent=2.6\zw\hangafter=1 $e^{(-1+2j)t}=e^{-t}e^{j2t}=e^{-t}(\cos2t+j\sin2t)$。実部は $e^{-t}\cos2t$\par
\bkvika $t=0$：$e^0\cos0=1\times1=1$\par
\bkvika $t=\dfrac{\pi}{4}$：$\cos\dfrac{\pi}{2}=0$ なので0\par
\bkvika $t=\dfrac{\pi}{2}$：$e^{-\pi/2}\cos\pi=0.208\times(-1)=-0.208$\par
\par\vspace{3pt}{\centering
\begin{tikzpicture}
\begin{axis}[bkplot, scale only axis, width=70mm, height=24mm, xmin=0, xmax=3.4, ymin=-0.45, ymax=1.1, xtick={0.785,1.571}, xticklabels={$\frac{\pi}{4}$,$\frac{\pi}{2}$}, ytick={1}, grid=none,
    axis x line=bottom, x axis line style={draw=none}, xticklabel style={yshift=-0.5mm}]
  \draw[bkInk, line width=0.8pt, -{Stealth[length=4pt]}] (axis cs:0,0) -- (axis cs:3.5,0) node[anchor=west, font=\sffamily\footnotesize, text=bkInk] {$t$};
  \addplot[bkPurple, line width=1.5pt, domain=0:3.3, samples=120] {exp(-x)*cos(deg(2*x))};
  \addplot[only marks, mark=*, mark size=1.6pt, bkInk] coordinates {(0,1) (0.785,0) (1.571,-0.208)};
\end{axis}
\end{tikzpicture}\par}\vspace{2pt}
\kotae\hangindent=2.6\zw\hangafter=1 $e^{-t}\cos2t$。1、0、$-0.208$
\end{reidai}

\begin{jibun}
$e^{(-0.5+3j)t}$ の実部を $t$ の式で書き、$t=0$、$\dfrac{\pi}{3}$ のときの値を求めましょう（$e^{-\pi/6}=0.592$）。
\bkjans{$e^{-0.5t}\cos3t$。$t=0$ で1。$t=\dfrac{\pi}{3}$ で $e^{-\pi/6}\cos\pi=-0.592$。}
\end{jibun}

\begin{reidai}[オイラーの公式で2倍角の式を作る]
$(e^{j\theta})^2=e^{j2\theta}$ を使って、$\cos2\theta$ と $\sin2\theta$ を、$\cos\theta$ と $\sin\theta$ で表しましょう。\par
\kangaekata 左は $(\cos\theta+j\sin\theta)^2$ を展開、右は $\cos2\theta+j\sin2\theta$。実部どうし、虚部どうしを等しいとおく（④と同じ）。\par
\shiki\hangindent=2.6\zw\hangafter=1 左：$(\cos\theta+j\sin\theta)^2=\cos^2\theta+2j\sin\theta\cos\theta+j^2\sin^2\theta$\par
\bkvikb $=(\cos^2\theta-\sin^2\theta)+j(2\sin\theta\cos\theta)$\par
\bkvika 右：$e^{j2\theta}=\cos2\theta+j\sin2\theta$\par
\bkvika 実部どうし・虚部どうしを等しいとおく\par
\par\vspace{3pt}{\centering
\begin{tikzpicture}[x=1mm, y=1mm]
  \draw[bkGray, line width=0.6pt, -{Stealth[length=3pt]}] (-4,0) -- (22,0) node[anchor=north east, font=\sffamily\footnotesize, inner sep=1pt, yshift=-1mm] {実部};
  \draw[bkGray, line width=0.6pt, -{Stealth[length=3pt]}] (0,-3) -- (0,20) node[anchor=north west, font=\sffamily\footnotesize, inner sep=1pt, xshift=1mm] {虚部};
  \draw[bkInk!40, line width=0.7pt] (16,0) arc (0:95:16);
  \draw[bkBlue, line width=1.2pt, -{Stealth[length=4pt]}] (0,0) -- (25:16);
  \draw[bkPurple, line width=1.2pt, -{Stealth[length=4pt]}] (0,0) -- (50:16);
  \node[font=\sffamily\footnotesize, inner sep=1pt, text=bkBlue, anchor=west] at ($(25:16)+(0.8,0)$) {$e^{j\theta}$};
  \node[font=\sffamily\footnotesize, inner sep=1pt, text=bkPurple, anchor=south west] at ($(50:16)+(0.5,0.3)$) {$(e^{j\theta})^2=e^{j2\theta}$（角が2倍）};
\end{tikzpicture}\par}\vspace{2pt}
\kotae\hangindent=2.6\zw\hangafter=1 $\cos2\theta=\cos^2\theta-\sin^2\theta$、$\sin2\theta=2\sin\theta\cos\theta$
\end{reidai}

\begin{jibun}
$e^{ja}e^{-jb}=e^{j(a-b)}$ を使って、$\cos(a-b)$ と $\sin(a-b)$ を、$\cos a$、$\sin a$、$\cos b$、$\sin b$ で表しましょう。
\bkjans{$(\cos a+j\sin a)(\cos b-j\sin b)$ を展開すると、実部は $\cos a\cos b+\sin a\sin b$、虚部は $\sin a\cos b-\cos a\sin b$。右は $\cos(a-b)+j\sin(a-b)$。よって $\cos(a-b)=\cos a\cos b+\sin a\sin b$、$\sin(a-b)=\sin a\cos b-\cos a\sin b$。}
\end{jibun}



%% ---------------- まちがえやすいところ・物理工学 ----------------
\bksection{tsumazuki}{まちがえやすいところをチェック}

\begin{tsumazuki}{角を度のまま肩に入れる}
\bkng{$e^{j60}=\dfrac{1}{2}+\dfrac{\sqrt3}{2}j$}
\bkok{$e^{j\pi/3}=\dfrac{1}{2}+\dfrac{\sqrt3}{2}j$（60°は $\dfrac{\pi}{3}$ ラジアン）}
$e^{j\theta}$ の $\theta$ はラジアンです。③で $(\cos t)'=-\sin t$ を使ったように、この公式は、ラジアンで書いた微分のきまり（5-15）から作ったからです。$e^{j60}$ は、60ラジアン（約3438°）回した点になってしまいます。
\end{tsumazuki}

\begin{tsumazuki}{$e^{j\theta}$ の大きさが、$\theta$ とともにふえると思う}
\bkng{$\lvert e^{j2}\rvert=e^2\approx7.39$}
\bkok{$\lvert e^{j2}\rvert=\sqrt{\cos^22+\sin^22}=1$}
$e^{j\theta}$ は、どんな $\theta$ でも単位円の上の点です（4-8 の $\sin^2\theta+\cos^2\theta=1$）。肩が実数の $e^{x}$ とちがい、肩に $j$ がつくと「ふえる」ではなく「回る」になります。
\end{tsumazuki}

\begin{tsumazuki}{$\sin$ の式で、$j$ でわるのを忘れる}
\bkng{$\sin\theta=\dfrac{e^{j\theta}-e^{-j\theta}}{2}$}
\bkok{$\sin\theta=\dfrac{e^{j\theta}-e^{-j\theta}}{2j}$}
ひくと $2j\sin\theta$ が残るので、$2j$ でわります。$\theta=\dfrac{\pi}{2}$ で確かめると、$\dfrac{j-(-j)}{2}=j$ となって、$\sin\dfrac{\pi}{2}=1$ と合いません。
\end{tsumazuki}

\begin{tsumazuki}{肩のたし算を、ふつうのたし算にする}
\bkng{$e^{j(a+b)}=e^{ja}+e^{jb}$}
\bkok{$e^{j(a+b)}=e^{ja}\times e^{jb}$}
肩（指数）がたし算のとき、本体はかけ算です（3-11 の $2^{3+2}=2^3\times2^2$ と同じ）。$a=b=0$ で確かめると、左は $e^0=1$、×の右は $1+1=2$ で合いません。
\end{tsumazuki}

\begin{tsumazuki}{矢印のまま答える}
\bkng{電圧は $100e^{j\omega t}$〔V〕}
\bkok{電圧は $100\sin\omega t$〔V〕（矢印のたての影）}
電圧計ではかれるのは実数です。矢印は計算の道具で、最後に決めたほうの影をとって答えます。
\end{tsumazuki}

\begin{butsuri}
\bkvikbox{家のコンセントと交流の計算}
\noindent\begin{minipage}[c]{0.66\linewidth}
\raggedright
日本の家のコンセントは、実効値100\,V（振幅は約141\,V）の交流です。周波数は、東日本が50\,Hz、西日本が60\,Hz です。モーター・照明・パソコンの電源などがつながると、電流と電圧の山の位置がずれます。電気の技術者は、電圧と電流をフェーザで表し、「矢印の図」（フェーザ図）をかいて、どれだけずれているか、どれだけの電力が使われるかを計算します。
\end{minipage}\hfill
\begin{minipage}[c]{0.3\linewidth}
\centering
\begin{tikzpicture}[x=1mm, y=1mm]
  \draw[bkInk!70, line width=0.9pt, fill=white, rounded corners=2mm] (0,0) rectangle (18,26);
  \draw[bkInk!70, line width=0.7pt, fill=bkPaper, rounded corners=1.5mm] (3,5) rectangle (15,21);
  \fill[bkInk] (6.2,10) rectangle (7.4,16);
  \fill[bkInk] (10.6,10) rectangle (11.8,16.8);
  \begin{scope}[xshift=25mm, yshift=4mm]
  \draw[bkGray, line width=0.5pt] (-1,0) -- (14,0);
  \draw[bkBlue, line width=1.3pt, -{Stealth[length=4pt]}] (0,0) -- (13,0);
  \draw[bkOrange, line width=1.3pt, -{Stealth[length=4pt]}] (0,0) -- (-30:10);
  \node[font=\sffamily\footnotesize, text=bkBlue, anchor=south, inner sep=1pt] at (7,0.6) {電圧};
  \node[font=\sffamily\footnotesize, text=bkOrange!80!black, anchor=north, inner sep=1pt] at (5,-6.4) {電流};
  \end{scope}
\end{tikzpicture}
\end{minipage}
\end{butsuri}

\begin{butsuri}
\bkvikbox{三相交流の発電機}
\noindent\begin{minipage}[c]{0.3\linewidth}
\centering
\begin{tikzpicture}[x=1mm, y=1mm]
  \draw[bkInk!60, line width=1pt] (0,0) circle (15);
  \foreach \a/\c in {90/bkBlue, 210/bkOrange, 330/bkPurple} {\begin{scope}[rotate=\a] \draw[\c, line width=0.9pt, fill=\c!15, rounded corners=0.8mm] (10,-3.5) rectangle (14,3.5); \end{scope}}
  \begin{scope}[rotate=20]
  \fill[bkRed!80] (-6,-2.2) rectangle (0,2.2);
  \fill[bkBlue!70] (0,-2.2) rectangle (6,2.2);
  \node[font=\sffamily\bfseries\footnotesize, text=white] at (-3,0) {N};
  \node[font=\sffamily\bfseries\footnotesize, text=white] at (3,0) {S};
  \end{scope}
  \draw[bkInk, line width=0.7pt, -{Stealth[length=3pt]}] (40:8.6) arc (40:100:8.6);
\end{tikzpicture}
\end{minipage}\hfill
\begin{minipage}[c]{0.66\linewidth}
\raggedright
発電所の発電機では、まん中の磁石（電磁石）が回り、まわりに120°ずつずらして置いた3組のコイルに、電圧が生まれます。3つの電圧は、フェーザで書くと120°ずつずれた3本の矢印で、その和は0です（⑨）。だから、3つの負荷（電気を使う機械）が同じなら、もどりの電線を共通にしても、そこには電流がほとんど流れません。少ない電線で、大きな電力を送れます。
\end{minipage}
\end{butsuri}

\begin{butsuri}
\bkvikbox{スマホの電波と、点の回転}
\noindent\begin{minipage}[c]{0.66\linewidth}
\raggedright
スマホや Wi-Fi の電波は、「どれだけ強く、どの角だけずらして送るか」で情報を運びます。送る波を複素数（矢印）で表すと、送れる記号は、複素平面の上の点になります。たとえば16個の点を使えば、1回に4ビット（$2^4=16$）の情報を送れます。受けとった側は、ずれて回ってしまった点に $e^{-j\theta}$ をかけて、角をもどしてから読みます。オイラーの公式は、ポケットの中でも毎秒たくさん使われています。
\end{minipage}\hfill
\begin{minipage}[c]{0.3\linewidth}
\centering
\begin{tikzpicture}[x=1mm, y=1mm]
  \draw[bkGray, line width=0.6pt, -{Stealth[length=3pt]}] (-15,0) -- (16,0);
  \draw[bkGray, line width=0.6pt, -{Stealth[length=3pt]}] (0,-15) -- (0,16);
  \foreach \i in {-3,-1,1,3} \foreach \k in {-3,-1,1,3} {\fill[bkBlue] ({3.4*\i},{3.4*\k}) circle (0.9);}
  \draw[bkOrange, line width=1.2pt, -{Stealth[length=3.5pt]}] (0,0) -- (10.2,3.4);
\end{tikzpicture}
\end{minipage}
\end{butsuri}



\begin{column}{複素数は「計算の道具」だけ？}
この章では、回る矢印（複素数）を、交流の計算を楽にする道具として使い、最後に影（実数）をとりました。電気・機械の工学では、これで正しい使い方です。ところが、原子や電子の世界を表す量子力学では、事情がちがいます。電子のふるまいを表す「波動関数」は、はじめから $e^{j\theta}$ をふくむ複素数の波で、その矢印の大きさの2乗が「そこに電子が見つかりやすさ」（確率）を表します。複素数そのものが、自然のしくみを表しているのです。くわしくは物理シリーズ『原子』で学びます。
\end{column}

%% ---------------- 章末問題 ----------------
\bkendproblems
\bkproblemgroup{基本}{4}

\bkproblem 次の数を $a+bj$ の形に直しましょう。\newline（1）$e^{j\pi/2}$\quad（2）$e^{j5\pi/6}$\quad（3）$2e^{-j\pi/3}$
\bkpans{（1）\nobreak $j$。（2）\nobreak 150°なので $-\dfrac{\sqrt3}{2}+\dfrac{1}{2}j$。（3）\nobreak $2\Bigl(\dfrac{1}{2}-\dfrac{\sqrt3}{2}j\Bigr)=1-\sqrt3\,j$。}

\bkproblem 次の数を $re^{j\theta}$ の形に直しましょう（$-\pi<\theta\leqq\pi$）。\newline（1）$1+j$\quad（2）$-\sqrt3+j$
\bkpans{（1）\nobreak $\sqrt2\,e^{j\pi/4}$。（2）\nobreak $r=2$、左上で150°。$2e^{j5\pi/6}$。}

\bkproblem 次の計算をし、$a+bj$ の形で答えましょう。\newline（1）$4e^{j\pi/4}\times0.5e^{j\pi/4}$\quad（2）$\dfrac{8e^{j\pi}}{2e^{j\pi/3}}$
\bkpans{（1）\nobreak $2e^{j\pi/2}=2j$。（2）\nobreak $4e^{j2\pi/3}=-2+2\sqrt3\,j$。}

\bkproblem $(1+\sqrt3\,j)^3$ を、極形式を使って計算しましょう。
\bkpans{$1+\sqrt3\,j=2e^{j\pi/3}$。$2^3e^{j\pi}=8\times(-1)=-8$。}

\bkproblem $\theta=\dfrac{\pi}{2}$ のとき、$\dfrac{e^{j\theta}+e^{-j\theta}}{2}$ と $\dfrac{e^{j\theta}-e^{-j\theta}}{2j}$ を計算し、$\cos\dfrac{\pi}{2}$、$\sin\dfrac{\pi}{2}$ と同じになることを確かめましょう。
\bkpans{$e^{j\pi/2}=j$、$e^{-j\pi/2}=-j$。$\dfrac{j+(-j)}{2}=0=\cos\dfrac{\pi}{2}$。$\dfrac{j-(-j)}{2j}=\dfrac{2j}{2j}=1=\sin\dfrac{\pi}{2}$。}

\bkproblem 静電容量 \qty{20}{\micro F} のコンデンサーに、$v=100\sin\omega t$〔V〕、周波数 \qty{50}{Hz} の電圧をかけます。電流の振幅と、電圧とのずれを求めましょう（$\pi=3.14$）。
\bkpans{$\omega=314$〔rad/s〕。$314\times20\times10^{-6}\times100\approx0.628$。振幅 約\qty{0.63}{A}、電流は電圧より90°進む。}

\bkproblem $e^{j3t}$ を、1回微分、2回微分、積分（積分定数は考えない）しましょう。それぞれ、もとの矢印を何倍し、何度回したものですか。
\bkpans{微分 $3je^{j3t}$（3倍、90°進む）。2回微分 $-9e^{j3t}$（9倍、180°）。積分 $-\dfrac{j}{3}e^{j3t}$（$\dfrac{1}{3}$ 倍、90°おくれる）。}

\bkproblem インダクタンス \qty{0.05}{H} のコイルに、$v=100\sin\omega t$〔V〕、周波数 \qty{50}{Hz} の電圧をかけます。電流の振幅と、電圧とのずれを求めましょう（$\pi=3.14$）。
\bkpans{$j\omega L=314\times0.05j=15.7j$〔Ω〕。$\dfrac{100}{15.7j}\approx-6.37j$。振幅 約\qty{6.4}{A}、電流は電圧より90°おくれる。}

\bkproblemgroup{チャレンジ}{5}

\bkproblem 三相交流のフェーザ $1$、$e^{j2\pi/3}$、$e^{j4\pi/3}$ を、それぞれ $a+bj$ に直してたし、和が0になることを確かめましょう。
\bkpans{$e^{j2\pi/3}=-\dfrac{1}{2}+\dfrac{\sqrt3}{2}j$、$e^{j4\pi/3}=-\dfrac{1}{2}-\dfrac{\sqrt3}{2}j$。実部 $1-\dfrac{1}{2}-\dfrac{1}{2}=0$、虚部 $\dfrac{\sqrt3}{2}-\dfrac{\sqrt3}{2}=0$。}

\bkproblem $e^{(-0.5+2j)t}$ の実部を $t$ の式で書き、$t=0$、$\dfrac{\pi}{4}$、$\dfrac{\pi}{2}$ のときの値を求めましょう（$e^{-\pi/4}=0.456$）。
\bkpans{$e^{-0.5t}\cos2t$。$t=0$ で1、$t=\dfrac{\pi}{4}$ で0、$t=\dfrac{\pi}{2}$ で $0.456\times(-1)=-0.456$。}

\bkproblem 次の和を、フェーザを使って1つの $\sin$ の形で表しましょう。\newline（1）$100\sin\omega t+50\cos\omega t$（$\tan^{-1}0.5=26.6°$）。\newline（2）$v_1=100\sin\omega t$〔V〕と $v_2=100\sin(\omega t-60°)$〔V〕の和（角がマイナスのフェーザは右回り）。
\bkpans{（1）\nobreak フェーザの和 $100+50j$。大きさ $\sqrt{10000+2500}=\sqrt{12500}\approx111.8$、角 $\tan^{-1}\dfrac{50}{100}=26.6°$。$111.8\sin(\omega t+26.6°)$。（2）\nobreak $100+100e^{-j\pi/3}=100+50-86.6j=150-86.6j$。大きさ $\sqrt{30000}\approx173$、右下の点で角 $-30°$。$173\sin(\omega t-30°)$〔V〕。}

\bkproblem $e^{ja}e^{jb}=e^{j(a+b)}$ の両辺をオイラーの公式で書いて広げ、$\cos(a+b)$ と $\sin(a+b)$ の加法定理を作りましょう。また、$a=b=\theta$ とすると、2倍角の式になることを確かめましょう。
\bkpans{左を展開すると、実部は $\cos a\cos b-\sin a\sin b$、虚部は $\sin a\cos b+\cos a\sin b$。右は $\cos(a+b)+j\sin(a+b)$。実部・虚部を等しいとおく。$a=b=\theta$ とすると、$\cos2\theta=\cos^2\theta-\sin^2\theta$、$\sin2\theta=2\sin\theta\cos\theta$。}

\begin{furikaeri}
\begin{checklist}
\item $e^{j\theta}$ が単位円の上の角 $\theta$ の点になるわけを、小さな回転のくり返しと微分で説明できる。\bkfurref{bk:6-10:kaiten}
\item $e^{j\theta}$ を $a+bj$ に直し、$re^{j\theta}$ のかけ算・わり算・べきを計算できる。\bkfurref{bk:6-10:kyoku}
\item フェーザのたし算で交流を合成し、微分が $j\omega$ をかけることだと説明できる。\bkfurref{bk:6-10:phasor}
\end{checklist}
\end{furikaeri}

\bkfukidashi{sora}{$e^{j\theta}$ は「回る」、$e^{x}$ は「ふえる」。次の章では、$\cos$ のなかまの「ふえる版」、双曲線関数に出会うよ。}
